\documentclass[a4paper,11pt]{article}
        \usepackage[top=15mm,bottom=18mm,left=16mm,right=16mm]{geometry}
        \usepackage{fontspec}
        \usepackage{xeCJK}
        \setmainfont{Times New Roman}
        \setCJKmainfont{Yu Gothic}
        \setmonofont{Consolas}
        \setCJKmonofont{Yu Gothic}
        \usepackage{booktabs}
        \usepackage{longtable}
        \usepackage{tabularx}
        \usepackage{array}
        \usepackage{enumitem}
        \usepackage{hyperref}
        \usepackage{listings}
        \usepackage{xcolor}
        \usepackage{amsmath}
        \usepackage{amssymb}
        \hypersetup{
          colorlinks=true,
          linkcolor=blue,
          urlcolor=blue,
          pdftitle={目的関数重み CEM 学習},
          pdfauthor={Codex}
        }
        \lstset{
          basicstyle=\ttfamily\small,
          breaklines=true,
          columns=fullflexible,
          keepspaces=true,
          frame=single,
          framerule=0.2pt,
          rulecolor=\color{black},
          showstringspaces=false
        }
        \setlength{\parskip}{0.42em}
        \setlength{\parindent}{1em}
        \setlength{\emergencystretch}{3em}
        \renewcommand{\arraystretch}{1.15}
        \title{36. 目的関数重み CEM 学習}
        \author{solar\_ws0129-main}
        \date{2026-07-29}
        \begin{document}
        \maketitle
        \tableofcontents
        \clearpage

        \section{対象}
        \begin{tabularx}{\linewidth}{>{\raggedright\arraybackslash}p{0.18\linewidth} X}
        \toprule
        項目 & 内容 \\
        \midrule
        ファイル & \path|scripts/tune_upper_planner_weights.py| \\
        ソースSHA-256 & \path|7fa117bb2a11983fb9101004a230ec22fb899a1c95d928fecb2e4743f7a0b22f| \\
        種別 & Python \\
        区分 & planning research \\
        役割 & 複数シナリオで upper planner の cost weight を探索し、risk-aware な candidate を評価する。 \\
        起動文脈 & learn action の中心。 \\
        \bottomrule
        \end{tabularx}

        \section{このファイルを読む理由}
        複数シナリオで upper planner の cost weight を探索し、risk-aware な candidate を評価する。

        \begin{itemize}[leftmargin=1.6em]
        \item multi-scenario、CVaR、chance gate を扱う。
\item operational release ではなく exact validation 前段の性格が強い。
        \end{itemize}

        \section{呼び出し関係}
        \subsection{どこから呼ばれるか}
        \begin{itemize}[leftmargin=1.6em]
        \item \path|scripts/solar_control.sh|
        \end{itemize}

        \subsection{次に読むと理解が進むファイル}
        \begin{itemize}[leftmargin=1.6em]
        \item \path|mpc_solarcar/upper_cost.py|
\item \path|scripts/solar_sim.py|
        \end{itemize}

        \subsection{同一リポジトリ内の主要依存}
        \begin{itemize}[leftmargin=1.6em]
        \item \path|mpc_solarcar/risk_utils.py|
\item \path|mpc_solarcar/schedule_utils.py|
\item \path|mpc_solarcar/upper_cost.py|
        \end{itemize}

        \section{ソース冒頭の把握}
        モジュール docstring または先頭意図の要約:

        モジュール docstring は明示されていない。

        \subsection{代表コード断片}
        \begin{lstlisting}[language=Python]
@dataclass
class TermSpec:
    name: str
    lo: float
    hi: float
    init_log10: float
    threshold: float = 1.0e-4


@dataclass
class ScenarioSpec:
    name: str
    cfg_overrides: Dict[str, object]
    cli_overrides: Dict[str, object]
    weight: float


def read_yaml(path: Path) -> dict:
    with path.open("r", encoding="utf-8") as f:
        return yaml.safe_load(f) or {}
\end{lstlisting}


        \section{主要構造の読み取り}
        主要クラスは TermSpec, ScenarioSpec。 主要関数は read\_yaml, write\_yaml, ensure\_dir, log\_trial\_to\_tensorboard, repo\_relative, failed\_scenario\_result, compile\_tex, latex\_escape。 CLI 引数宣言は 19 件。

        \subsection{主要クラス・関数}
            \begin{longtable}{p{0.07\linewidth} p{0.16\linewidth} p{0.25\linewidth} p{0.40\linewidth}}
            \toprule
            行 & 種別 & 名前 & 補足 \\
            \midrule
            83 & class & \path|TermSpec| &  \\
92 & class & \path|ScenarioSpec| &  \\
99 & function & \path|read_yaml| & args: path \\
104 & function & \path|write_yaml| & args: path, payload \\
110 & function & \path|ensure_dir| & args: path \\
114 & function & \path|log_trial_to_tensorboard| & args: writer, prefix, result, step \\
147 & function & \path|repo_relative| & args: path\_like \\
159 & function & \path|failed_scenario_result| & args: scenario\_name, scenario\_weight, upper\_cost \\
217 & function & \path|compile_tex| & args: tex\_path \\
234 & function & \path|latex_escape| & args: text \\
253 & function & \path|format_override_value| & args: value \\
265 & function & \path|canonical_runtime_weights| & args: weights \\
270 & function & \path|set_nested_value| & args: cfg, dotted\_key, value \\
283 & function & \path|speed_series| & args: df \\
291 & function & \path|upper_cost_specs| & args: cfg \\
336 & function & \path|vector_to_weights| & args: specs, vec, base\_cfg \\
347 & function & \path|mirror_legacy_weights| & args: cfg, upper\_cost \\
360 & function & \path|build_reference_free_profile| & args: profile\_yaml, output\_dir \\
397 & function & \path|default_scenarios| & args: profile\_yaml \\
482 & function & \path|run_single_scenario| & args: profile\_yaml, output\_dir, candidate\_name, scenario, upper\_cost, cfg\_overrides, cli\_overrides \\
613 & function & \path|evaluate_simulation| & args: profile\_yaml, summary, sim\_df, detail\_df, upper\_cost \\
759 & function & \path|aggregate_candidate| & args: candidate, scenario\_results, weights, risk\_config \\
820 & function & \path|run_candidate| & args: profile\_yaml, output\_dir, candidate\_name, upper\_cost, cfg\_overrides, cli\_overrides, scenarios, risk\_config \\
858 & function & \path|save_trial_checkpoint| & args: output\_dir, trials, best\_result \\
868 & function & \path|csv_row_count| & args: path \\
873 & function & \path|summarize_path| & args: path \\
896 & function & \path|dataframe_to_markdown| & args: df \\
910 & function & \path|flatten_scalars| & args: prefix, payload, rows \\
920 & function & \path|build_current_asset_manifests| & args: profile\_yaml, fit\_summary, output\_dir \\
986 & function & \path|render_report| & args: output\_dir, source\_profile\_yaml, eval\_profile\_yaml, fit\_summary, baseline\_result, tuned\_result, trials\_df, best\_weights, manifest\_paths, removed\_reference, specs, search\_cfg, tensorboard\_dir \\
1138 & function & \path|tex_path_rel| & args: path \\
1504 & function & \path|main| &  \\
            \bottomrule
            \end{longtable}

        \subsection{ファイルを上から読んだときの定義順}
            \begin{longtable}{p{0.07\linewidth} p{0.84\linewidth}}
            \toprule
            行 & 内容 \\
            \midrule
            19 & matplotlib.use(...) を実行する。 \\
25 & 例外処理を伴う try ブロックを実行する。 \\
30 & ROOT に Path(\_\_file\_\_).resolve().parents[1] の結果を代入する。 \\
31 & 条件 str(ROOT) not in sys.path を判定し、真なら内部処理を行う。 \\
39 & DEFAULT\_PROFILE に ROOT / 'project\_packages' / 'bwsc2025\_fitted\_mle4' / 'profile.yaml' の結果を代入する。 \\
42 & LITERATURE に [\{'label': 'de Boer et al. (2005)', 'title': 'A Tutorial on the Cross-Entropy Method', 'url': 'https://doi.org/10.1007/s10479-005-5724-z', 'note': 'CEM is a generic derivative-free method for hard optimization problems.'\}, \{'label': 'Gros and Zanon (2019/2020)', 'title': 'Data-driven Economic NMPC using Reinforcement Learning', 'url': 'https://arxiv.org/pdf/1904.04152', 'note': 'RL can tune stage cost, terminal cost, and constraints of MPC/Economic MPC.'\}, \{'label': 'Zarrouki et al. (2024)', 'title': 'A Safe Reinforcement Learning driven Weights-varying Model Predictive Controller', 'url': 'https://arxiv.org/pdf/2402.02624', 'note': 'Safe RL can adapt MPC weights within a restricted safe search space.'\}, \{'label': 'Howlett et al. (1997)', 'title': 'Optimal driving strategy for a solar car on a level road', 'url': 'https://doi.org/10.1093/imaman/8.1.59', 'note': 'Solar-race strategy is governed by a tight energy-speed trade-off.'\}, \{'label': 'Pudney and Howlett (2002)', 'title': 'Critical Speed Control of a Solar Car', 'url': 'https://link.springer.com/article/10.1023/A\%3A1020907101234', 'note': 'Large unnecessary speed deviations are undesirable in solar-race operation.'\}, \{'label': 'Byrd et al. (1995)', 'title': 'A Limited Memory Algorithm for Bound Constrained Optimization', 'url': 'https://doi.org/10.1137/0916069', 'note': 'L-BFGS-B provides bounded local refinement after global candidate search.'\}] の結果を代入する。 \\
83 & クラス TermSpec を定義する。 \\
92 & クラス ScenarioSpec を定義する。 \\
99 & 関数 read\_yaml を定義する。 \\
104 & 関数 write\_yaml を定義する。 \\
110 & 関数 ensure\_dir を定義する。 \\
114 & 関数 log\_trial\_to\_tensorboard を定義する。 \\
147 & 関数 repo\_relative を定義する。 \\
159 & 関数 failed\_scenario\_result を定義する。 \\
217 & 関数 compile\_tex を定義する。 \\
234 & 関数 latex\_escape を定義する。 \\
253 & 関数 format\_override\_value を定義する。 \\
265 & 関数 canonical\_runtime\_weights を定義する。 \\
270 & 関数 set\_nested\_value を定義する。 \\
283 & 関数 speed\_series を定義する。 \\
291 & 関数 upper\_cost\_specs を定義する。 \\
336 & 関数 vector\_to\_weights を定義する。 \\
347 & 関数 mirror\_legacy\_weights を定義する。 \\
360 & 関数 build\_reference\_free\_profile を定義する。 \\
397 & 関数 default\_scenarios を定義する。 \\
482 & 関数 run\_single\_scenario を定義する。 \\
613 & 関数 evaluate\_simulation を定義する。 \\
759 & 関数 aggregate\_candidate を定義する。 \\
820 & 関数 run\_candidate を定義する。 \\
858 & 関数 save\_trial\_checkpoint を定義する。 \\
868 & 関数 csv\_row\_count を定義する。 \\
873 & 関数 summarize\_path を定義する。 \\
896 & 関数 dataframe\_to\_markdown を定義する。 \\
910 & 関数 flatten\_scalars を定義する。 \\
920 & 関数 build\_current\_asset\_manifests を定義する。 \\
986 & 関数 render\_report を定義する。 \\
1504 & 関数 main を定義する。 \\
1940 & 条件 \_\_name\_\_ == '\_\_main\_\_' を判定し、真なら内部処理を行う。 \\
            \bottomrule
            \end{longtable}

        \subsection{import 群の詳細}
            \begin{longtable}{p{0.07\linewidth} p{0.33\linewidth} p{0.50\linewidth}}
            \toprule
            行 & import 文 & このファイルで読む理由 \\
            \midrule
            2 & \path|from __future__ import annotations| & 型ヒントの遅延評価を有効にし、自己参照型や forward reference を安全に扱うため。 このファイル内での使用位置は少ないか、間接利用である。 \\
4 & \path|import argparse| & CLI 引数を宣言し、実行時パラメータを外から受け取るため。 このファイル内での主な使用位置は L1505。 \\
5 & \path|import json| & manifest、checkpoint、UDP payload をやり取りするため。 このファイル内での主な使用位置は L571, L607, L852, L1936, L1937。 \\
6 & \path|import math| & Python標準のスカラー数値関数と定数を使うため。実際に参照する名前と行は使用位置欄で確認する。 このファイル内での主な使用位置は L259, L299, L300, L301, L302, L303, L304, L305, ...。 \\
7 & \path|import os| & パス、環境変数、プロセス外部状態を扱うため。 このファイル内での主な使用位置は L154, L205, L206, L207, L208, L209, L210, L211, ...。 \\
8 & \path|import shutil| & shutil モジュールを利用するため。 このファイル内での主な使用位置は L1862。 \\
9 & \path|import subprocess| & git 情報取得や外部コマンド実行を行うため。 このファイル内での主な使用位置は L220, L224, L225, L228, L545, L550。 \\
10 & \path|import sys| & sys モジュールを利用するため。 このファイル内での主な使用位置は L31, L32, L504。 \\
11 & \path|import textwrap| & textwrap モジュールを利用するため。 このファイル内での主な使用位置は L1444。 \\
12 & \path|from dataclasses import dataclass| & dataclasses から dataclass を読み込み、このファイルの処理を組み立てるため。 このファイル内での主な使用位置は L82, L91。 \\
13 & \path|from datetime import datetime| & UTC 時刻や相対時間を扱うため。 このファイル内での主な使用位置は L965, L1163, L1563。 \\
14 & \path|from pathlib import Path| & ファイルやディレクトリを安全に扱うため。 このファイル内での主な使用位置は L30, L99, L104, L110, L151, L167, L168, L169, ...。 \\
15 & \path|from typing import Dict, Iterable, List| & typing から Dict, Iterable, List を読み込み、このファイルの処理を組み立てるため。 このファイル内での主な使用位置は L94, L95, L114, L162, L165, L166, L174, L265, ...。 \\
17 & \path|import matplotlib| & matplotlib モジュールを利用するため。 このファイル内での主な使用位置は L19。 \\
20 & \path|import matplotlib.pyplot as plt| & matplotlib.pyplot モジュールを利用するため。 このファイル内での主な使用位置は L1037, L1038, L1039, L1040, L1041, L1042, L1043, L1044, ...。 \\
21 & \path|import numpy as np| & 数値配列、clip、補間、統計計算を行うため。 このファイル内での主な使用位置は L288, L336, L339, L627, L637, L639, L643, L653, ...。 \\
22 & \path|import pandas as pd| & CSV や時刻列の読込・整形・再サンプリングに使うため。 このファイル内での主な使用位置は L283, L288, L572, L573, L616, L617, L623, L626, ...。 \\
23 & \path|import yaml| & profile、stop、schedule、summary YAML を読み書きするため。 このファイル内での主な使用位置は L101, L107。 \\
26 & \path|from torch.utils.tensorboard import SummaryWriter| & torch.utils.tensorboard から SummaryWriter を読み込み、このファイルの処理を組み立てるため。 このファイル内での主な使用位置は L28, L1567。 \\
34 & \path|from mpc_solarcar.schedule_utils import DriveSchedule| & schedule\_utils.py から DriveSchedule を読み込み、このファイルの内部処理を分担させるため。 実体ファイルは mpc\_solarcar/schedule\_utils.py。 このファイル内での主な使用位置は L635。 \\
35 & \path|from mpc_solarcar.risk_utils import ScenarioRiskConfig, aggregate_scenario_scores| & risk\_utils.py から ScenarioRiskConfig, aggregate\_scenario\_scores を読み込み、このファイルの内部処理を分担させるため。 実体ファイルは mpc\_solarcar/risk\_utils.py。 このファイル内での主な使用位置は L763, L775, L828, L1536。 \\
36 & \path|from mpc_solarcar.upper_cost import UpperCostConfig, active_upper_cost_terms, load_upper_cost_config| & 上位MPC 目的関数 から UpperCostConfig, active\_upper\_cost\_terms, load\_upper\_cost\_config を読み込み、このファイルの内部処理を分担させるため。 実体ファイルは mpc\_solarcar/upper\_cost.py。 このファイル内での主な使用位置は L175, L292, L336, L684, L1135, L1534, L1552。 \\
            \bottomrule
            \end{longtable}

        \section{このファイルを読む前に必要な基礎知識}
        次の章は、構文やROS用語を既知と仮定しないための説明である。

        \subsection{プログラム、プロセス、メモリ、オブジェクトを区別する}
ソースファイルはディスク上の文字列であり、それ自体は走っていない。Python実行可能プログラムがソースを読み、OSがその実行を一つのプロセスとして管理する。プロセスには仮想メモリ、開いているファイル、スレッド、終了コードなどが対応する。


\begin{lstlisting}
ソースファイル -> Pythonインタプリタが読み込む -> OS上のプロセス -> Pythonオブジェクトをメモリに生成 -> 関数やcallbackを実行
\end{lstlisting}


「メモリ上に生成する」とは、実行中プロセスが使う記憶領域に、その値の型、属性、参照関係を表す実体を用意することである。変数は実体そのものというより、そのオブジェクトを指す名前として理解するとPythonの代入が読みやすい。


\begin{lstlisting}[language=python]
node = MPCNode()
alias = node
# nodeとaliasは同じオブジェクトを参照する。
\end{lstlisting}


一つのプロセスには複数スレッドを持たせられる。スレッドは同じプロセスのメモリを共有するため、受信callbackと最適化callbackが同じself.zを書き換える場合は実行順と排他を検討する必要がある。

\paragraph{根拠資料}
\begin{itemize}[leftmargin=1.6em]
\item \href{https://docs.python.org/3/tutorial/classes.html}{Python公式チュートリアル: Classes}
\end{itemize}

\subsection{名前、代入、型、参照}
Pythonの代入文は、右辺を先に評価し、その結果のオブジェクトへ左辺の名前を結び付ける。`x = f()`では、まずfを呼び、その戻り値が得られてからxが更新される。


`float(...)`、`int(...)`、`bool(...)`は型名を呼び出して値を変換する。外部CSV、YAML、ROS parameterから来た値は型が期待どおりとは限らないため、このリポジトリでは明示変換が多い。


`None`は値が無いことを示す単一のオブジェクト、`math.nan`は浮動小数点値ではあるが有効な数値ではないことを示す。両者は用途が異なるため、`is None`と`math.isfinite`を使い分ける。

\paragraph{根拠資料}
\begin{itemize}[leftmargin=1.6em]
\item \href{https://docs.python.org/3/tutorial/classes.html}{Python公式チュートリアル: Classes}
\end{itemize}

\subsection{関数、引数、戻り値、スコープ}
`def`は関数本体をその場で実行する文ではない。関数オブジェクトを作り、その名前を現在の名前空間へ登録する。`f`は関数そのもの、`f()`はその関数を呼んだ結果である。


仮引数は関数定義側の受け取り名、実引数は呼出側から渡す値である。位置引数、キーワード引数、既定値付き引数、`*args`、`**kwargs`は渡し方が異なる。


\begin{lstlisting}[language=python]
def clip_speed(value, minimum=0.0, maximum=110.0):
    return min(maximum, max(minimum, value))

v = clip_speed(120.0, maximum=90.0)
\end{lstlisting}


関数内で作った通常の名前はローカルスコープに属する。メソッドが`self.z`を更新すればオブジェクトの状態が残るが、単に`z`を更新しただけならその呼出しのローカル値である。

\paragraph{根拠資料}
\begin{itemize}[leftmargin=1.6em]
\item \href{https://docs.python.org/3/tutorial/classes.html}{Python公式チュートリアル: Classes}
\end{itemize}

\subsection{module、package、import、console\_scripts}
Pythonファイルはmoduleとして読み込める。複数moduleをディレクトリにまとめたものがpackageである。`from .model import SolarCarModel`の先頭の点は、現在と同じpackage内のmodel moduleを指す相対importである。


import時にはファイルのトップレベル文が上から一度実行される。`def`や`class`は関数・クラスを登録するが、その本体の通常処理は呼び出すまで走らない。


setuptoolsの`console\_scripts`は、端末で使う実行可能名と`package.module:function`を対応付けるインストール時メタデータである。生成される実行用ラッパーはmoduleをimportし、指定関数を呼び、戻り値を終了コードとして扱う小さな入口である。


\begin{lstlisting}
ROS launchのexecutable名 -> install済みconsole script -> mpc_solarcar.mpc_nodeをimport -> main()を呼ぶ
\end{lstlisting}

\paragraph{根拠資料}
\begin{itemize}[leftmargin=1.6em]
\item \href{https://setuptools.pypa.io/en/latest/userguide/entry_point.html}{setuptools公式: Entry Points / Console Scripts}
\item \href{https://docs.python.org/3/tutorial/classes.html}{Python公式チュートリアル: Classes}
\end{itemize}

\subsection{例外、try/finally、with、資源解放}
例外は通常の戻り値とは別経路で異常を呼出元へ伝える。`try/except`は想定した異常を処理し、`finally`は成功・失敗にかかわらず後始末を行う。


`with open(...) as f:`はcontext managerを使い、ブロックを出るとファイルを閉じる。CSVやログの破損を避けるため、開いた資源の所有者と閉じる場所を明確にする。


`except Exception: pass`はノードを止めない利点がある一方、入力異常を隠して原因追跡を難しくする。安全に関係する値では、少なくとも頻度制限付き警告、異常カウンタ、fallback状態のpublishを検討する。



\subsection{class、独自クラス、インスタンス、self、継承}
クラスは、保持するデータと、そのデータを扱う関数を一つの型としてまとめる設計単位である。「独自クラス」とはPythonやROS 2が最初から用意した型ではなく、このプロジェクトが目的に合わせて定義した新しい型を指す。


`class MPCNode(Node)`は、rclpyのNodeを基底クラスとして継承し、Nodeが持つpublisher、subscription、timer、parameterなどの機能にMPC固有機能を追加する。丸括弧内は関数引数ではなく基底クラス指定である。


`MPCNode()`はクラスオブジェクトを呼び出して新しいインスタンスを作る。Pythonは新しい実体を用意し、その実体を第1引数selfとして`\_\_init\_\_`へ渡す。`self`は予約語ではなく慣習的な名前だが、この慣習を守ることでコードの意味が共有される。


\begin{lstlisting}[language=python]
class Counter:
    def __init__(self):
        self.value = 0

    def increment(self):
        self.value += 1

counter = Counter()
counter.increment()
\end{lstlisting}


`super().\_\_init\_\_(...)`は継承元の初期化処理を呼ぶ。MPCNodeではこれを省くとROSノードとして必要な内部実体が作られず、`create\_publisher`などを正しく使えない。

\paragraph{根拠資料}
\begin{itemize}[leftmargin=1.6em]
\item \href{https://docs.python.org/3/tutorial/classes.html}{Python公式チュートリアル: Classes}
\item \href{https://docs.ros.org/en/humble/Tutorials/Beginner-CLI-Tools/Understanding-ROS2-Nodes/Understanding-ROS2-Nodes.html}{ROS 2 Humble公式: Understanding nodes}
\end{itemize}

\subsection{先頭アンダースコア、\_\_init\_\_、名前付けの作法}
`self.\_step\_solar`の先頭1個のアンダースコアは「クラス内部で使う実装詳細」という慣習であり、アクセスを強制禁止する機構ではない。外部コードから呼べるが、公開APIとして依存しない意思表示である。


`\_\_init\_\_`の前後2個のアンダースコアはPythonが定めた特殊メソッド名である。クラス生成時に自動で呼ばれる。任意の名前に前後2個を付けて独自仕様を作ることは避ける。


`\_`で始まるローカル変数は、未使用または外部へ見せない意図を示す場合がある。例えば`\_base\_csv`は戻り値として受け取る必要はあるが、その後の処理では使わないことを読者へ示す。

\paragraph{根拠資料}
\begin{itemize}[leftmargin=1.6em]
\item \href{https://docs.python.org/3/tutorial/classes.html}{Python公式チュートリアル: Classes}
\end{itemize}

\subsection{デコレータと@dataclass}
`@名前`は、直後に定義した関数またはクラスを別の関数へ渡し、その結果で定義名を置き換えるデコレータ構文である。`@Class`という一般構文があるのではなく、`@dataclass`など実際のデコレータ名を書く。


`@dataclass`は型注釈付きフィールドから`\_\_init\_\_`、`\_\_repr\_\_`、比較処理などを自動生成する。物理パラメータや解析結果のように、名前付きデータを一まとまりで運ぶ用途に向く。


\begin{lstlisting}[language=python]
from dataclasses import dataclass

@dataclass
class State:
    soc: float
    temperature_c: float

state = State(soc=0.8, temperature_c=30.0)
\end{lstlisting}


自動生成は検証を自動で保証しない。単位、許容範囲、相互依存制約は`\_\_post\_init\_\_`や別の検証関数で確認する必要がある。

\paragraph{根拠資料}
\begin{itemize}[leftmargin=1.6em]
\item \href{https://docs.python.org/3/library/dataclasses.html}{Python公式ライブラリ: dataclasses}
\item \href{https://docs.python.org/3/tutorial/classes.html}{Python公式チュートリアル: Classes}
\end{itemize}

\subsection{list、dict、deque、NumPy配列、pandas DataFrame}
listは順序付き可変列、tupleは順序付きで通常変更しない列、dictはキーから値を引く対応表、dequeは両端追加・削除が効率的なキューである。どの構造を選ぶかはアクセス方法と更新方法で決まる。


NumPy配列は同種数値を連続的に扱い、要素ごとの演算、clip、補間、線形代数を簡潔に書く。shapeは各次元の要素数であり、速度系列なら通常`(N,)`、候補集団なら`(population, N)`となる。


pandas DataFrameは列名を持つ表である。CSV読込後は列型、欠損、単位、timezone、並び順、重複を明示的に処理しなければ、数値計算が動いても意味が誤る。



\subsection{CLI、PowerShell、Bash、環境変数、終了コード}
CLIは端末からプログラム名と引数を渡す操作界面である。`argparse`は文字列として届く引数を名前、型、既定値、必須性に従って解析する。


PowerShellとBashは別のshellであり、変数記法、改行継続、引用、パス表記が異なる。このプロジェクトではWindows側のSolarSim.ps1がWSL側のsolar\_control.shへ処理を渡す。


環境変数は親プロセスから子プロセスへ受け渡される名前付き文字列である。ROS\_DOMAIN\_ID、RMW\_IMPLEMENTATION、Pythonの数値スレッド数などはコード外から動作を変えるため、実行記録へ残す必要がある。


終了コード0は一般に成功、0以外は失敗を示す。shellルータは子プロセスの終了コードを握り潰さず上位へ返すことで、自動運用が失敗を検知できる。

\paragraph{根拠資料}
\begin{itemize}[leftmargin=1.6em]
\item \href{https://setuptools.pypa.io/en/latest/userguide/entry_point.html}{setuptools公式: Entry Points / Console Scripts}
\end{itemize}

\subsection{目的関数、制約、L-BFGS-B、SHGO、有限grid証明}
数値最適化器は、利用者が与えた目的関数を複数の候補点で評価し、より小さい値を持つ候補を探す。solverが物理を理解するのではなく、物理と運用価値はcost関数へ書かれる。


L-BFGS-Bは変数ごとの上下限を扱える局所最適化法である。初期値の近くの谷へ収束し得るため、非凸問題では複数seedや大域探索と組み合わせる。successがFalseでも有限な候補が返る場合があるため、採用条件をコード側で決める。


SHGOは定めたsamplingと局所最適化を組み合わせる大域最適化法である。有限Cartesian gridの全列挙は、そのgrid上の最良を証明できるが、連続領域全体の最良を自動的に証明しない。資料ではこの証明範囲を区別する。

\paragraph{根拠資料}
\begin{itemize}[leftmargin=1.6em]
\item \href{https://docs.scipy.org/doc/scipy/reference/generated/scipy.optimize.minimize.html}{SciPy公式: scipy.optimize.minimize}
\item \href{https://docs.scipy.org/doc/scipy/reference/generated/scipy.optimize.shgo.html}{SciPy公式: scipy.optimize.shgo}
\end{itemize}

\subsection{Cross-Entropy Methodを式と実装で理解する}
CEMは候補を生成する確率分布を持ち、良かったelite候補から分布を更新する反復的な確率最適化である。このリポジトリのupper\_solver.pyは各制御点速度を独立正規分布で生成し、上下限へclipする。


\[
u_i^{(j)} \sim \mathcal{N}(\mu_i^{(g)},(\sigma_i^{(g)})^2),\qquad u_i^{(j)}\leftarrow\operatorname{clip}(u_i^{(j)},l_i,h_i)
\]

\[
\mu_i^{(g+1)}=\frac{1}{K}\sum_{j\in\mathcal{E}_g}u_i^{(j)},\qquad \sigma_i^{(g+1)}=\max\left(\operatorname{Std}_{j\in\mathcal{E}_g}u_i^{(j)},0.05(h_i-l_i)\right)
\]

ここでE\_gはcostが小さい上位K候補である。平均は良い領域へ移り、標準偏差は探索幅を表す。標準偏差の下限は探索が完全に潰れることを避ける。


現行hybrid\_bounded\_minimizeは、deterministic seedを評価し、上位候補をL-BFGS-Bで局所refineし、設定とseed間不一致に応じてCEMを実行し、最後に再度局所refineする。したがってCEM単独ではなくhybrid solverである。


CEMで落とした候補を永久保存しないこと自体は通常の最適化として自然だが、off-nominal状態からの再利用には別のpolicy library設計が必要である。状態を無制限に全組合せ保存する代わりに、SoC、進捗、時刻、予報誤差、停止状態などのscenarioを設計し、近傍policyを検索してMPCで再最適化する。

\paragraph{根拠資料}
\begin{itemize}[leftmargin=1.6em]
\item \href{https://link.springer.com/book/10.1007/978-1-4757-4321-0}{Rubinstein and Kroese: The Cross-Entropy Method}
\item \href{https://docs.scipy.org/doc/scipy/reference/generated/scipy.optimize.minimize.html}{SciPy公式: scipy.optimize.minimize}
\end{itemize}

\subsection{制御、状態、入力、モデル予測制御MPC}
制御対象の内部を表す状態をx、操作入力をu、外乱・予報をwとすると、離散モデルは`x[k+1] = f(x[k], u[k], w[k])`と書ける。ソーラーカーではSoC、電池温度、距離、速度などが状態候補、速度目標や駆動トルクが入力候補になる。


\[
\min_{u_0,\ldots,u_{N-1}} \sum_{k=0}^{N-1}\ell(x_k,u_k,w_k)+V_f(x_N)
\]

MPCは現在状態からNステップ先まで予測し、目的関数と制約を満たす入力系列を求める。ただし実際に適用するのは通常先頭入力だけで、次回は新しい実測状態から再び解く。これがreceding horizonである。


予測モデル、目的関数、制約、ホライズン、solver、初期値のどれかが変わると答えも変わる。「MPCを使う」だけでは仕様は決まらず、これらを単位付きで追う必要がある。

\paragraph{根拠資料}
\begin{itemize}[leftmargin=1.6em]
\item \href{https://docs.scipy.org/doc/scipy/reference/generated/scipy.optimize.minimize.html}{SciPy公式: scipy.optimize.minimize}
\end{itemize}



        \section{関数・クラスを上から順に解説}
        \subsection{L83 クラス TermSpec}
            \begin{itemize}[leftmargin=1.6em]
              \item 行範囲: L83--L88
              \item 所属: module直下
              \item 定義: \path|TermSpec(bases=none)|
              \item docstring: 明示されていない。
              \item このブロックが直接呼ぶ主な関数・メソッド: 特に明示されていない。
              \item 戻り値の要点: 戻り値が無いか、return 文が明示されていない。
              \item 主なローカル代入名: 特になし。
              \item 読み取る主なインスタンス属性: 特になし。
              \item 更新する主なインスタンス属性: 特になし。
              \item 明示的に送出する例外: 明示的raiseはない。
              \item 制御構造の規模: 条件分岐 0、ループ 0、try 0
            \end{itemize}
            \paragraph{この定義を読むためのPython構文}
            \begin{itemize}[leftmargin=1.6em]
            \item class文は新しい型を定義する。丸括弧内は継承する基底クラスである。
\item @で始まる行は、定義したクラスを加工するdecoratorである。
            \end{itemize}
            \paragraph{上から順の処理}
            \begin{enumerate}[leftmargin=1.8em]
            \item name に  を代入する。
\item lo に  を代入する。
\item hi に  を代入する。
\item init\_log10 に  を代入する。
\item threshold に 0.0001 を代入する。
            \end{enumerate}
            \paragraph{代表コード断片}
            \begin{lstlisting}[language=Python]
class TermSpec:
    name: str
    lo: float
    hi: float
    init_log10: float
    threshold: float = 1.0e-4
\end{lstlisting}

\subsection{L92 クラス ScenarioSpec}
            \begin{itemize}[leftmargin=1.6em]
              \item 行範囲: L92--L96
              \item 所属: module直下
              \item 定義: \path|ScenarioSpec(bases=none)|
              \item docstring: 明示されていない。
              \item このブロックが直接呼ぶ主な関数・メソッド: 特に明示されていない。
              \item 戻り値の要点: 戻り値が無いか、return 文が明示されていない。
              \item 主なローカル代入名: 特になし。
              \item 読み取る主なインスタンス属性: 特になし。
              \item 更新する主なインスタンス属性: 特になし。
              \item 明示的に送出する例外: 明示的raiseはない。
              \item 制御構造の規模: 条件分岐 0、ループ 0、try 0
            \end{itemize}
            \paragraph{この定義を読むためのPython構文}
            \begin{itemize}[leftmargin=1.6em]
            \item class文は新しい型を定義する。丸括弧内は継承する基底クラスである。
\item @で始まる行は、定義したクラスを加工するdecoratorである。
            \end{itemize}
            \paragraph{上から順の処理}
            \begin{enumerate}[leftmargin=1.8em]
            \item name に  を代入する。
\item cfg\_overrides に  を代入する。
\item cli\_overrides に  を代入する。
\item weight に  を代入する。
            \end{enumerate}
            \paragraph{代表コード断片}
            \begin{lstlisting}[language=Python]
class ScenarioSpec:
    name: str
    cfg_overrides: Dict[str, object]
    cli_overrides: Dict[str, object]
    weight: float
\end{lstlisting}

\subsection{L99 関数 read\_yaml}
            \begin{itemize}[leftmargin=1.6em]
              \item 行範囲: L99--L101
              \item 所属: module直下
              \item 定義: \path|read_yaml(path: Path) -> dict|
              \item docstring: 明示されていない。
              \item このブロックが直接呼ぶ主な関数・メソッド: open, safe\_load
              \item 戻り値の要点: yaml.safe\_load(f) or \{\}
              \item 主なローカル代入名: f
              \item 読み取る主なインスタンス属性: 特になし。
              \item 更新する主なインスタンス属性: 特になし。
              \item 明示的に送出する例外: 明示的raiseはない。
              \item 制御構造の規模: 条件分岐 0、ループ 0、try 0
            \end{itemize}
            \paragraph{この定義を読むためのPython構文}
            \begin{itemize}[leftmargin=1.6em]
            \item `->`以後は戻り値の型注釈であり、実行時の値を自動変換する命令ではない。
\item with文はcontext managerに開始・終了処理を任せ、ファイルなどの資源を確実に解放する。
            \end{itemize}
            \paragraph{上から順の処理}
            \begin{enumerate}[leftmargin=1.8em]
            \item with 文で path.open('r', encoding='utf-8') を管理しながら処理する。
\item   yaml.safe\_load(f) or \{\} を返す。
            \end{enumerate}
            \paragraph{代表コード断片}
            \begin{lstlisting}[language=Python]
def read_yaml(path: Path) -> dict:
    with path.open("r", encoding="utf-8") as f:
        return yaml.safe_load(f) or {}
\end{lstlisting}

\subsection{L104 関数 write\_yaml}
            \begin{itemize}[leftmargin=1.6em]
              \item 行範囲: L104--L107
              \item 所属: module直下
              \item 定義: \path|write_yaml(path: Path, payload: dict) -> None|
              \item docstring: 明示されていない。
              \item このブロックが直接呼ぶ主な関数・メソッド: mkdir, open, safe\_dump
              \item 戻り値の要点: 戻り値が無いか、return 文が明示されていない。
              \item 主なローカル代入名: f
              \item 読み取る主なインスタンス属性: 特になし。
              \item 更新する主なインスタンス属性: 特になし。
              \item 明示的に送出する例外: 明示的raiseはない。
              \item 制御構造の規模: 条件分岐 0、ループ 0、try 0
            \end{itemize}
            \paragraph{この定義を読むためのPython構文}
            \begin{itemize}[leftmargin=1.6em]
            \item `->`以後は戻り値の型注釈であり、実行時の値を自動変換する命令ではない。
\item with文はcontext managerに開始・終了処理を任せ、ファイルなどの資源を確実に解放する。
            \end{itemize}
            \paragraph{上から順の処理}
            \begin{enumerate}[leftmargin=1.8em]
            \item path.parent.mkdir(...) を実行する。
\item with 文で path.open('w', encoding='utf-8', newline='\textbackslash{}n') を管理しながら処理する。
\item   yaml.safe\_dump(...) を実行する。
            \end{enumerate}
            \paragraph{代表コード断片}
            \begin{lstlisting}[language=Python]
def write_yaml(path: Path, payload: dict) -> None:
    path.parent.mkdir(parents=True, exist_ok=True)
    with path.open("w", encoding="utf-8", newline="\n") as f:
        yaml.safe_dump(payload, f, sort_keys=False, allow_unicode=True)
\end{lstlisting}

\subsection{L110 関数 ensure\_dir}
            \begin{itemize}[leftmargin=1.6em]
              \item 行範囲: L110--L111
              \item 所属: module直下
              \item 定義: \path|ensure_dir(path: Path) -> None|
              \item docstring: 明示されていない。
              \item このブロックが直接呼ぶ主な関数・メソッド: mkdir
              \item 戻り値の要点: 戻り値が無いか、return 文が明示されていない。
              \item 主なローカル代入名: 特になし。
              \item 読み取る主なインスタンス属性: 特になし。
              \item 更新する主なインスタンス属性: 特になし。
              \item 明示的に送出する例外: 明示的raiseはない。
              \item 制御構造の規模: 条件分岐 0、ループ 0、try 0
            \end{itemize}
            \paragraph{この定義を読むためのPython構文}
            \begin{itemize}[leftmargin=1.6em]
            \item `->`以後は戻り値の型注釈であり、実行時の値を自動変換する命令ではない。
            \end{itemize}
            \paragraph{上から順の処理}
            \begin{enumerate}[leftmargin=1.8em]
            \item path.mkdir(...) を実行する。
            \end{enumerate}
            \paragraph{代表コード断片}
            \begin{lstlisting}[language=Python]
def ensure_dir(path: Path) -> None:
    path.mkdir(parents=True, exist_ok=True)
\end{lstlisting}

\subsection{L114 関数 log\_trial\_to\_tensorboard}
            \begin{itemize}[leftmargin=1.6em]
              \item 行範囲: L114--L144
              \item 所属: module直下
              \item 定義: \path|log_trial_to_tensorboard(writer, prefix: str, result: Dict[str, object], step: int) -> None|
              \item docstring: 明示されていない。
              \item このブロックが直接呼ぶ主な関数・メソッド: add\_scalar, float
              \item 戻り値の要点: 戻り値が無いか、return 文が明示されていない。
              \item 主なローカル代入名: key, scalar\_keys
              \item 読み取る主なインスタンス属性: 特になし。
              \item 更新する主なインスタンス属性: 特になし。
              \item 明示的に送出する例外: 明示的raiseはない。
              \item 制御構造の規模: 条件分岐 1、ループ 1、try 1
            \end{itemize}
            \paragraph{この定義を読むためのPython構文}
            \begin{itemize}[leftmargin=1.6em]
            \item `->`以後は戻り値の型注釈であり、実行時の値を自動変換する命令ではない。
\item try文は例外が起き得る処理と、異常時または終了時の経路を分ける。
            \end{itemize}
            \paragraph{上から順の処理}
            \begin{enumerate}[leftmargin=1.8em]
            \item 条件 writer is None を判定し、真なら内部処理を行う。
\item    を返す。
\item scalar\_keys に ['score', 'score\_mean', 'score\_worst', 'score\_variance', 'loss\_cvar', 'scenario\_failure\_probability', 'chance\_constraint\_pass', 'final\_distance\_km', 'final\_distance\_worst\_km', 'avg\_speed\_kmh', 'min\_soc', 'final\_soc', 'oscillation\_mean\_abs\_dv\_kmh', 'oscillation\_p95\_abs\_dv\_kmh', 'current\_rms\_a', 'pack\_slew\_rms\_kw', 'daylight\_stop\_h', 'daylight\_full\_soc\_h', 'unused\_finish\_soc', 'cpu\_sec', 'active\_term\_count'] の結果を代入する。
\item scalar\_keys を順に走査し、各要素を key に入れて処理する。
\item   例外処理を伴う try ブロックを実行する。
\item     writer.add\_scalar(...) を実行する。
\item     Exceptionを捕捉した場合:
\item     Continue 文を実行する。
            \end{enumerate}
            \paragraph{代表コード断片}
            \begin{lstlisting}[language=Python]
def log_trial_to_tensorboard(writer, prefix: str, result: Dict[str, object], step: int) -> None:
    if writer is None:
        return
    scalar_keys = [
        "score",
        "score_mean",
        "score_worst",
        "score_variance",
        "loss_cvar",
        "scenario_failure_probability",
        "chance_constraint_pass",
        "final_distance_km",
        "final_distance_worst_km",
        "avg_speed_kmh",
        "min_soc",
        "final_soc",
        "oscillation_mean_abs_dv_kmh",
        "oscillation_p95_abs_dv_kmh",
        "current_rms_a",
        "pack_slew_rms_kw",
        "daylight_stop_h",
        "daylight_full_soc_h",
        "unused_finish_soc",
        "cpu_sec",
        "active_term_count",
    ]
    for key in scalar_keys:
        try:
            writer.add_scalar(f"{prefix}/{key}", float(result[key]), step)
        except Exception:
            continue
\end{lstlisting}

\subsection{L147 関数 repo\_relative}
            \begin{itemize}[leftmargin=1.6em]
              \item 行範囲: L147--L156
              \item 所属: module直下
              \item 定義: \path|repo_relative(path_like) -> str|
              \item docstring: 明示されていない。
              \item このブロックが直接呼ぶ主な関数・メソッド: Path, fspath, is\_absolute, relative\_to, replace, resolve, str, strip
              \item 戻り値の要点: '' / os.fspath(resolved.relative\_to(ROOT)).replace('\textbackslash{}\textbackslash{}', '/') / raw.replace('\textbackslash{}\textbackslash{}', '/')
              \item 主なローカル代入名: path, raw, resolved
              \item 読み取る主なインスタンス属性: 特になし。
              \item 更新する主なインスタンス属性: 特になし。
              \item 明示的に送出する例外: 明示的raiseはない。
              \item 制御構造の規模: 条件分岐 1、ループ 0、try 1
            \end{itemize}
            \paragraph{この定義を読むためのPython構文}
            \begin{itemize}[leftmargin=1.6em]
            \item `->`以後は戻り値の型注釈であり、実行時の値を自動変換する命令ではない。
\item try文は例外が起き得る処理と、異常時または終了時の経路を分ける。
            \end{itemize}
            \paragraph{上から順の処理}
            \begin{enumerate}[leftmargin=1.8em]
            \item raw に str(path\_like or '').strip() の結果を代入する。
\item 条件 not raw を判定し、真なら内部処理を行う。
\item   '' を返す。
\item path に Path(raw) の結果を代入する。
\item 例外処理を伴う try ブロックを実行する。
\item   resolved に path.resolve() if path.is\_absolute() else (ROOT / path).resolve() の結果を代入する。
\item   os.fspath(resolved.relative\_to(ROOT)).replace('\textbackslash{}\textbackslash{}', '/') を返す。
\item   Exceptionを捕捉した場合:
\item   raw.replace('\textbackslash{}\textbackslash{}', '/') を返す。
            \end{enumerate}
            \paragraph{代表コード断片}
            \begin{lstlisting}[language=Python]
def repo_relative(path_like) -> str:
    raw = str(path_like or "").strip()
    if not raw:
        return ""
    path = Path(raw)
    try:
        resolved = path.resolve() if path.is_absolute() else (ROOT / path).resolve()
        return os.fspath(resolved.relative_to(ROOT)).replace("\\", "/")
    except Exception:
        return raw.replace("\\", "/")
\end{lstlisting}

\subsection{L159 関数 failed\_scenario\_result}
            \begin{itemize}[leftmargin=1.6em]
              \item 行範囲: L159--L214
              \item 所属: module直下
              \item 定義: \path|failed_scenario_result(scenario_name: str, scenario_weight: float, upper_cost: Dict[str, float], *, error: str, cfg_overrides: Dict[str, object], cli_overrides: Dict[str, object], summary_json: Path, out_csv: Path, detail_csv: Path, plan_csv: Path, report_html: Path, resolved_yaml: Path, sim_log_path: Path) -> Dict[str, object]|
              \item docstring: 明示されていない。
              \item このブロックが直接呼ぶ主な関数・メソッド: UpperCostConfig, active\_upper\_cost\_terms, float, fspath, int, len
              \item 戻り値の要点: \{'score': -1000000000.0, 'final\_distance\_km': 0.0, 'avg\_speed\_kmh': 0.0, 'min\_soc': 0.0, 'final\_soc': 0.0, 'elapsed\_hours': 0.0, 'cpu\_sec': 0.0, 'finish\_reached': False, 'model\_validation\_gate\_pass': False, 'scenario\_feasible': False, 'scenario\_feasibility\_checks': \{'simulation\_completed': False\}, 'scenario\_infeasibility\_reasons': [error], 'oscillation\_mean\_abs\_dv\_kmh': 1000000.0, 'oscillation\_p95\_abs\_dv\_kmh': 1000000.0, 'current\_rms\_a': 1000000.0, 'pack\_slew\_rms\_kw': 1000000.0, 'high\_speed\_h': 1000000.0, 'daylight\_stop\_h': 1000000.0, 'daylight\_full\_soc\_h': 1000000.0, 'unused\_finish\_soc': 1.0, 'finish\_soc\_target': 0.0, 'terminal\_soc\_error': 1.0, 'active\_term\_count': int(len(active\_terms)), 'active\_terms': active\_terms, 'scenario': scenario\_name, 'scenario\_weight': float(scenario\_weight), 'cfg\_overrides': cfg\_overrides, 'cli\_overrides': cli\_overrides, 'summary\_json': os.fspath(summary\_json), 'out\_csv': os.fspath(out\_csv), 'detail\_csv': os.fspath(detail\_csv), 'plan\_csv': os.fspath(plan\_csv), 'report\_html': os.fspath(report\_html), 'resolved\_yaml': os.fspath(resolved\_yaml), 'simulation\_log': os.fspath(sim\_log\_path), 'failed': True, 'error': error\}
              \item 主なローカル代入名: active\_terms
              \item 読み取る主なインスタンス属性: 特になし。
              \item 更新する主なインスタンス属性: 特になし。
              \item 明示的に送出する例外: 明示的raiseはない。
              \item 制御構造の規模: 条件分岐 0、ループ 0、try 0
            \end{itemize}
            \paragraph{この定義を読むためのPython構文}
            \begin{itemize}[leftmargin=1.6em]
            \item `->`以後は戻り値の型注釈であり、実行時の値を自動変換する命令ではない。
            \end{itemize}
            \paragraph{上から順の処理}
            \begin{enumerate}[leftmargin=1.8em]
            \item active\_terms に active\_upper\_cost\_terms(UpperCostConfig(**upper\_cost), threshold=0.0001) の結果を代入する。
\item \{'score': -1000000000.0, 'final\_distance\_km': 0.0, 'avg\_speed\_kmh': 0.0, 'min\_soc': 0.0, 'final\_soc': 0.0, 'elapsed\_hours': 0.0, 'cpu\_sec': 0.0, 'finish\_reached': False, 'model\_validation\_gate\_pass': False, 'scenario\_feasible': False, 'scenario\_feasibility\_checks': \{'simulation\_completed': False\}, 'scenario\_infeasibility\_reasons': [error], 'oscillation\_mean\_abs\_dv\_kmh': 1000000.0, 'oscillation\_p95\_abs\_dv\_kmh': 1000000.0, 'current\_rms\_a': 1000000.0, 'pack\_slew\_rms\_kw': 1000000.0, 'high\_speed\_h': 1000000.0, 'daylight\_stop\_h': 1000000.0, 'daylight\_full\_soc\_h': 1000000.0, 'unused\_finish\_soc': 1.0, 'finish\_soc\_target': 0.0, 'terminal\_soc\_error': 1.0, 'active\_term\_count': int(len(active\_terms)), 'active\_terms': active\_terms, 'scenario': scenario\_name, 'scenario\_weight': float(scenario\_weight), 'cfg\_overrides': cfg\_overrides, 'cli\_overrides': cli\_overrides, 'summary\_json': os.fspath(summary\_json), 'out\_csv': os.fspath(out\_csv), 'detail\_csv': os.fspath(detail\_csv), 'plan\_csv': os.fspath(plan\_csv), 'report\_html': os.fspath(report\_html), 'resolved\_yaml': os.fspath(resolved\_yaml), 'simulation\_log': os.fspath(sim\_log\_path), 'failed': True, 'error': error\} を返す。
            \end{enumerate}
            \paragraph{代表コード断片}
            \begin{lstlisting}[language=Python]
def failed_scenario_result(
    scenario_name: str,
    scenario_weight: float,
    upper_cost: Dict[str, float],
    *,
    error: str,
    cfg_overrides: Dict[str, object],
    cli_overrides: Dict[str, object],
    summary_json: Path,
    out_csv: Path,
    detail_csv: Path,
    plan_csv: Path,
    report_html: Path,
    resolved_yaml: Path,
    sim_log_path: Path,
) -> Dict[str, object]:
    active_terms = active_upper_cost_terms(UpperCostConfig(**upper_cost), threshold=1.0e-4)
    return {
        "score": -1.0e9,
        "final_distance_km": 0.0,
        "avg_speed_kmh": 0.0,
        "min_soc": 0.0,
        "final_soc": 0.0,
        "elapsed_hours": 0.0,
        "cpu_sec": 0.0,
        "finish_reached": False,
        "model_validation_gate_pass": False,
        "scenario_feasible": False,
        "scenario_feasibility_checks": {"simulation_completed": False},
        "scenario_infeasibility_reasons": [error],
        "oscillation_mean_abs_dv_kmh": 1.0e6,
        "oscillation_p95_abs_dv_kmh": 1.0e6,
        "current_rms_a": 1.0e6,
        "pack_slew_rms_kw": 1.0e6,
        "high_speed_h": 1.0e6,
...
\end{lstlisting}

\subsection{L217 関数 compile\_tex}
            \begin{itemize}[leftmargin=1.6em]
              \item 行範囲: L217--L231
              \item 所属: module直下
              \item 定義: \path|compile_tex(tex_path: Path) -> Path|
              \item docstring: 明示されていない。
              \item このブロックが直接呼ぶ主な関数・メソッド: CalledProcessError, FileNotFoundError, exists, range, run, with\_suffix
              \item 戻り値の要点: pdf\_path
              \item 主なローカル代入名: \_, pdf\_path, res
              \item 読み取る主なインスタンス属性: 特になし。
              \item 更新する主なインスタンス属性: 特になし。
              \item 明示的に送出する例外: FileNotFoundError(pdf\_path), subprocess.CalledProcessError(res.returncode, res.args)
              \item 制御構造の規模: 条件分岐 2、ループ 1、try 0
            \end{itemize}
            \paragraph{この定義を読むためのPython構文}
            \begin{itemize}[leftmargin=1.6em]
            \item `->`以後は戻り値の型注釈であり、実行時の値を自動変換する命令ではない。
            \end{itemize}
            \paragraph{上から順の処理}
            \begin{enumerate}[leftmargin=1.8em]
            \item pdf\_path に tex\_path.with\_suffix('.pdf') の結果を代入する。
\item range(2) を順に走査し、各要素を \_ に入れて処理する。
\item   res に subprocess.run(['xelatex', '-interaction=nonstopmode', tex\_path.name], cwd=tex\_path.parent, check=False, stdout=subprocess.DEVNULL, stderr=subprocess.DEVNULL) の結果を代入する。
\item   条件 res.returncode != 0 and (not pdf\_path.exists()) を判定し、真なら内部処理を行う。
\item     subprocess.CalledProcessError(res.returncode, res.args) を送出する。
\item 条件 not pdf\_path.exists() を判定し、真なら内部処理を行う。
\item   FileNotFoundError(pdf\_path) を送出する。
\item pdf\_path を返す。
            \end{enumerate}
            \paragraph{代表コード断片}
            \begin{lstlisting}[language=Python]
def compile_tex(tex_path: Path) -> Path:
    pdf_path = tex_path.with_suffix(".pdf")
    for _ in range(2):
        res = subprocess.run(
            ["xelatex", "-interaction=nonstopmode", tex_path.name],
            cwd=tex_path.parent,
            check=False,
            stdout=subprocess.DEVNULL,
            stderr=subprocess.DEVNULL,
        )
        if res.returncode != 0 and not pdf_path.exists():
            raise subprocess.CalledProcessError(res.returncode, res.args)
    if not pdf_path.exists():
        raise FileNotFoundError(pdf_path)
    return pdf_path
\end{lstlisting}

\subsection{L234 関数 latex\_escape}
            \begin{itemize}[leftmargin=1.6em]
              \item 行範囲: L234--L250
              \item 所属: module直下
              \item 定義: \path|latex_escape(text: str) -> str|
              \item docstring: 明示されていない。
              \item このブロックが直接呼ぶ主な関数・メソッド: items, replace, str
              \item 戻り値の要点: out
              \item 主なローカル代入名: dst, out, repl, src
              \item 読み取る主なインスタンス属性: 特になし。
              \item 更新する主なインスタンス属性: 特になし。
              \item 明示的に送出する例外: 明示的raiseはない。
              \item 制御構造の規模: 条件分岐 0、ループ 1、try 0
            \end{itemize}
            \paragraph{この定義を読むためのPython構文}
            \begin{itemize}[leftmargin=1.6em]
            \item `->`以後は戻り値の型注釈であり、実行時の値を自動変換する命令ではない。
            \end{itemize}
            \paragraph{上から順の処理}
            \begin{enumerate}[leftmargin=1.8em]
            \item repl に \{'\textbackslash{}\textbackslash{}': '\textbackslash{}\textbackslash{}textbackslash\{\}', '\&': '\textbackslash{}\textbackslash{}\&', '\%': '\textbackslash{}\textbackslash{}\%', '\$': '\textbackslash{}\textbackslash{}\$', '\#': '\textbackslash{}\textbackslash{}\#', '\_': '\textbackslash{}\textbackslash{}\_', '\{': '\textbackslash{}\textbackslash{}\{', '\}': '\textbackslash{}\textbackslash{}\}', '\textasciitilde{}': '\textbackslash{}\textbackslash{}textasciitilde\{\}', '\textasciicircum{}': '\textbackslash{}\textbackslash{}textasciicircum\{\}'\} の結果を代入する。
\item out に str(text) の結果を代入する。
\item repl.items() を順に走査し、各要素を (src, dst) に入れて処理する。
\item   out に out.replace(src, dst) の結果を代入する。
\item out を返す。
            \end{enumerate}
            \paragraph{代表コード断片}
            \begin{lstlisting}[language=Python]
def latex_escape(text: str) -> str:
    repl = {
        "\\": r"\textbackslash{}",
        "&": r"\&",
        "%": r"\%",
        "$": r"\$",
        "#": r"\#",
        "_": r"\_",
        "{": r"\{",
        "}": r"\}",
        "~": r"\textasciitilde{}",
        "^": r"\textasciicircum{}",
    }
    out = str(text)
    for src, dst in repl.items():
        out = out.replace(src, dst)
    return out
\end{lstlisting}

\subsection{L253 関数 format\_override\_value}
            \begin{itemize}[leftmargin=1.6em]
              \item 行範囲: L253--L262
              \item 所属: module直下
              \item 定義: \path|format_override_value(value) -> str|
              \item docstring: 明示されていない。
              \item このブロックが直接呼ぶ主な関数・メソッド: isfinite, isinstance, str
              \item 戻り値の要点: str(value) / 'true' if value else 'false' / str(value) / '0'
              \item 主なローカル代入名: 特になし。
              \item 読み取る主なインスタンス属性: 特になし。
              \item 更新する主なインスタンス属性: 特になし。
              \item 明示的に送出する例外: 明示的raiseはない。
              \item 制御構造の規模: 条件分岐 4、ループ 0、try 0
            \end{itemize}
            \paragraph{この定義を読むためのPython構文}
            \begin{itemize}[leftmargin=1.6em]
            \item `->`以後は戻り値の型注釈であり、実行時の値を自動変換する命令ではない。
            \end{itemize}
            \paragraph{上から順の処理}
            \begin{enumerate}[leftmargin=1.8em]
            \item 条件 isinstance(value, bool) を判定し、真なら内部処理を行う。
\item   'true' if value else 'false' を返す。
\item 条件 isinstance(value, int) を判定し、真なら内部処理を行う。
\item   str(value) を返す。
\item 条件 isinstance(value, float) を判定し、真なら内部処理を行う。
\item   条件 math.isfinite(value) を判定し、真なら内部処理を行う。
\item     f'\{value:.12g\}' を返す。
\item   '0' を返す。
\item str(value) を返す。
            \end{enumerate}
            \paragraph{代表コード断片}
            \begin{lstlisting}[language=Python]
def format_override_value(value) -> str:
    if isinstance(value, bool):
        return "true" if value else "false"
    if isinstance(value, int):
        return str(value)
    if isinstance(value, float):
        if math.isfinite(value):
            return f"{value:.12g}"
        return "0"
    return str(value)
\end{lstlisting}

\subsection{L265 関数 canonical\_runtime\_weights}
            \begin{itemize}[leftmargin=1.6em]
              \item 行範囲: L265--L267
              \item 所属: module直下
              \item 定義: \path|canonical_runtime_weights(weights: Dict[str, float]) -> Dict[str, float]|
              \item docstring: Return the exact numeric values serialized into solar\_sim overrides.
              \item このブロックが直接呼ぶ主な関数・メソッド: float, format\_override\_value, items, str
              \item 戻り値の要点: \{str(key): float(format\_override\_value(float(value))) for key, value in weights.items()\}
              \item 主なローカル代入名: key, value
              \item 読み取る主なインスタンス属性: 特になし。
              \item 更新する主なインスタンス属性: 特になし。
              \item 明示的に送出する例外: 明示的raiseはない。
              \item 制御構造の規模: 条件分岐 0、ループ 0、try 0
            \end{itemize}
            \paragraph{この定義を読むためのPython構文}
            \begin{itemize}[leftmargin=1.6em]
            \item `->`以後は戻り値の型注釈であり、実行時の値を自動変換する命令ではない。
\item 内包表記は、反復・条件・値生成を一つの式で記述してcollectionを作る。
            \end{itemize}
            \paragraph{上から順の処理}
            \begin{enumerate}[leftmargin=1.8em]
            \item \{str(key): float(format\_override\_value(float(value))) for key, value in weights.items()\} を返す。
            \end{enumerate}
            \paragraph{代表コード断片}
            \begin{lstlisting}[language=Python]
def canonical_runtime_weights(weights: Dict[str, float]) -> Dict[str, float]:
    """Return the exact numeric values serialized into solar_sim overrides."""
    return {str(key): float(format_override_value(float(value))) for key, value in weights.items()}
\end{lstlisting}

\subsection{L270 関数 set\_nested\_value}
            \begin{itemize}[leftmargin=1.6em]
              \item 行範囲: L270--L280
              \item 所属: module直下
              \item 定義: \path|set_nested_value(cfg: dict, dotted_key: str, value) -> None|
              \item docstring: 明示されていない。
              \item このブロックが直接呼ぶ主な関数・メソッド: get, isinstance, split, str
              \item 戻り値の要点: 戻り値が無いか、return 文が明示されていない。
              \item 主なローカル代入名: child, part, parts, target
              \item 読み取る主なインスタンス属性: 特になし。
              \item 更新する主なインスタンス属性: 特になし。
              \item 明示的に送出する例外: 明示的raiseはない。
              \item 制御構造の規模: 条件分岐 2、ループ 1、try 0
            \end{itemize}
            \paragraph{この定義を読むためのPython構文}
            \begin{itemize}[leftmargin=1.6em]
            \item `->`以後は戻り値の型注釈であり、実行時の値を自動変換する命令ではない。
\item 内包表記は、反復・条件・値生成を一つの式で記述してcollectionを作る。
            \end{itemize}
            \paragraph{上から順の処理}
            \begin{enumerate}[leftmargin=1.8em]
            \item target に cfg の結果を代入する。
\item parts に [part for part in str(dotted\_key).split('.') if part] の結果を代入する。
\item parts[:-1] を順に走査し、各要素を part に入れて処理する。
\item   child に target.get(part) の結果を代入する。
\item   条件 not isinstance(child, dict) を判定し、真なら内部処理を行う。
\item     child に \{\} の結果を代入する。
\item     target[part] に child の結果を代入する。
\item   target に child の結果を代入する。
\item 条件 parts を判定し、真なら内部処理を行う。
\item   target[parts[-1]] に value の結果を代入する。
            \end{enumerate}
            \paragraph{代表コード断片}
            \begin{lstlisting}[language=Python]
def set_nested_value(cfg: dict, dotted_key: str, value) -> None:
    target = cfg
    parts = [part for part in str(dotted_key).split(".") if part]
    for part in parts[:-1]:
        child = target.get(part)
        if not isinstance(child, dict):
            child = {}
            target[part] = child
        target = child
    if parts:
        target[parts[-1]] = value
\end{lstlisting}

\subsection{L283 関数 speed\_series}
            \begin{itemize}[leftmargin=1.6em]
              \item 行範囲: L283--L288
              \item 所属: module直下
              \item 定義: \path|speed_series(df: pd.DataFrame) -> pd.Series|
              \item docstring: 明示されていない。
              \item このブロックが直接呼ぶ主な関数・メソッド: Series, astype, len, zeros
              \item 戻り値の要点: pd.Series(np.zeros(len(df), dtype=float)) / df['v\_exec\_kmh'].astype(float) / df['v\_cmd\_kmh'].astype(float)
              \item 主なローカル代入名: 特になし。
              \item 読み取る主なインスタンス属性: 特になし。
              \item 更新する主なインスタンス属性: 特になし。
              \item 明示的に送出する例外: 明示的raiseはない。
              \item 制御構造の規模: 条件分岐 2、ループ 0、try 0
            \end{itemize}
            \paragraph{この定義を読むためのPython構文}
            \begin{itemize}[leftmargin=1.6em]
            \item `->`以後は戻り値の型注釈であり、実行時の値を自動変換する命令ではない。
            \end{itemize}
            \paragraph{上から順の処理}
            \begin{enumerate}[leftmargin=1.8em]
            \item 条件 'v\_exec\_kmh' in df.columns を判定し、真なら内部処理を行う。
\item   df['v\_exec\_kmh'].astype(float) を返す。
\item 条件 'v\_cmd\_kmh' in df.columns を判定し、真なら内部処理を行う。
\item   df['v\_cmd\_kmh'].astype(float) を返す。
\item pd.Series(np.zeros(len(df), dtype=float)) を返す。
            \end{enumerate}
            \paragraph{代表コード断片}
            \begin{lstlisting}[language=Python]
def speed_series(df: pd.DataFrame) -> pd.Series:
    if "v_exec_kmh" in df.columns:
        return df["v_exec_kmh"].astype(float)
    if "v_cmd_kmh" in df.columns:
        return df["v_cmd_kmh"].astype(float)
    return pd.Series(np.zeros(len(df), dtype=float))
\end{lstlisting}

\subsection{L291 関数 upper\_cost\_specs}
            \begin{itemize}[leftmargin=1.6em]
              \item 行範囲: L291--L333
              \item 所属: module直下
              \item 定義: \path|upper_cost_specs(cfg: UpperCostConfig, *, include_progress_terms: bool = False, include_uncertainty_term: bool = True, include_terminal_term: bool = True) -> List[TermSpec]|
              \item docstring: 明示されていない。
              \item このブロックが直接呼ぶ主な関数・メソッド: TermSpec, append, extend, log10, max
              \item 戻り値の要点: specs
              \item 主なローカル代入名: specs
              \item 読み取る主なインスタンス属性: 特になし。
              \item 更新する主なインスタンス属性: 特になし。
              \item 明示的に送出する例外: 明示的raiseはない。
              \item 制御構造の規模: 条件分岐 3、ループ 0、try 0
            \end{itemize}
            \paragraph{この定義を読むためのPython構文}
            \begin{itemize}[leftmargin=1.6em]
            \item `->`以後は戻り値の型注釈であり、実行時の値を自動変換する命令ではない。
            \end{itemize}
            \paragraph{上から順の処理}
            \begin{enumerate}[leftmargin=1.8em]
            \item specs に [TermSpec('w\_speed\_smooth', -2.0, 3.0, math.log10(max(cfg.w\_speed\_smooth, 1e-06))), TermSpec('w\_speed\_limit', -2.0, 3.0, math.log10(max(cfg.w\_speed\_limit, 1e-06))), TermSpec('w\_current\_sq', -5.0, 1.0, math.log10(max(cfg.w\_current\_sq, 1e-06))), TermSpec('w\_pack\_energy', -5.0, 2.0, math.log10(max(cfg.w\_pack\_energy, 1e-06))), TermSpec('w\_joule\_loss', -5.0, 2.0, math.log10(max(cfg.w\_joule\_loss, 1e-06))), TermSpec('w\_aero\_energy', -5.0, 2.0, math.log10(max(cfg.w\_aero\_energy, 1e-06))), TermSpec('w\_mech\_energy', -5.0, 2.0, math.log10(max(cfg.w\_mech\_energy, 1e-06))), TermSpec('w\_kinetic\_pos', -5.0, 2.0, math.log10(max(cfg.w\_kinetic\_pos, 1e-06))), TermSpec('w\_pack\_power\_slew', -4.0, 4.0, math.log10(max(cfg.w\_pack\_power\_slew, 1e-06))), TermSpec('w\_speed\_quartic', -6.0, 2.0, math.log10(max(cfg.w\_speed\_quartic, 1e-06))), TermSpec('w\_solar\_headroom', -2.0, 7.0, max(2.0, math.log10(max(cfg.w\_solar\_headroom, 1e-06)))), TermSpec('w\_soc\_floor\_barrier', -6.0, 4.0, math.log10(max(cfg.w\_soc\_floor\_barrier, 1e-06))), TermSpec('w\_terminal\_soc\_min', -1.0, 4.0, math.log10(max(cfg.w\_terminal\_soc\_min, 1e-06))), TermSpec('w\_day\_end\_soc\_min', 2.0, 6.0, math.log10(max(cfg.w\_day\_end\_soc\_min, 1e-06)))] の結果を代入する。
\item 条件 include\_uncertainty\_term を判定し、真なら内部処理を行う。
\item   specs.append(...) を実行する。
\item 条件 include\_terminal\_term を判定し、真なら内部処理を行う。
\item   specs.append(...) を実行する。
\item 条件 include\_progress\_terms を判定し、真なら内部処理を行う。
\item   specs.extend(...) を実行する。
\item specs を返す。
            \end{enumerate}
            \paragraph{代表コード断片}
            \begin{lstlisting}[language=Python]
def upper_cost_specs(
    cfg: UpperCostConfig,
    *,
    include_progress_terms: bool = False,
    include_uncertainty_term: bool = True,
    include_terminal_term: bool = True,
) -> List[TermSpec]:
    specs = [
        TermSpec("w_speed_smooth", -2.0, 3.0, math.log10(max(cfg.w_speed_smooth, 1.0e-6))),
        TermSpec("w_speed_limit", -2.0, 3.0, math.log10(max(cfg.w_speed_limit, 1.0e-6))),
        TermSpec("w_current_sq", -5.0, 1.0, math.log10(max(cfg.w_current_sq, 1.0e-6))),
        TermSpec("w_pack_energy", -5.0, 2.0, math.log10(max(cfg.w_pack_energy, 1.0e-6))),
        TermSpec("w_joule_loss", -5.0, 2.0, math.log10(max(cfg.w_joule_loss, 1.0e-6))),
        TermSpec("w_aero_energy", -5.0, 2.0, math.log10(max(cfg.w_aero_energy, 1.0e-6))),
        TermSpec("w_mech_energy", -5.0, 2.0, math.log10(max(cfg.w_mech_energy, 1.0e-6))),
        TermSpec("w_kinetic_pos", -5.0, 2.0, math.log10(max(cfg.w_kinetic_pos, 1.0e-6))),
        TermSpec("w_pack_power_slew", -4.0, 4.0, math.log10(max(cfg.w_pack_power_slew, 1.0e-6))),
        TermSpec("w_speed_quartic", -6.0, 2.0, math.log10(max(cfg.w_speed_quartic, 1.0e-6))),
        # At high SoC the headroom term is multiplied by a small squared SoC
        # excess.  Its useful scale is therefore much larger than the legacy
        # 1e-4 seed; allow CEM to explore a physically meaningful range.
        TermSpec("w_solar_headroom", -2.0, 7.0, max(2.0, math.log10(max(cfg.w_solar_headroom, 1.0e-6)))),
        TermSpec("w_soc_floor_barrier", -6.0, 4.0, math.log10(max(cfg.w_soc_floor_barrier, 1.0e-6))),
        TermSpec("w_terminal_soc_min", -1.0, 4.0, math.log10(max(cfg.w_terminal_soc_min, 1.0e-6))),
        TermSpec("w_day_end_soc_min", 2.0, 6.0, math.log10(max(cfg.w_day_end_soc_min, 1.0e-6))),
    ]
    if include_uncertainty_term:
        specs.append(TermSpec("w_uncertainty_reserve", -6.0, 5.0, math.log10(max(cfg.w_uncertainty_reserve, 1.0e-6))))
    if include_terminal_term:
        specs.append(TermSpec("w_soc_terminal", -6.0, 6.0, math.log10(max(cfg.w_soc_terminal, 1.0e-6))))
    if include_progress_terms:
        specs.extend(
            [
                TermSpec("w_progress_lag", -6.0, 4.0, math.log10(max(cfg.w_progress_lag, 1.0e-6))),
                TermSpec(
...
\end{lstlisting}

\subsection{L336 関数 vector\_to\_weights}
            \begin{itemize}[leftmargin=1.6em]
              \item 行範囲: L336--L344
              \item 所属: module直下
              \item 定義: \path|vector_to_weights(specs: List[TermSpec], vec: np.ndarray, base_cfg: UpperCostConfig) -> Dict[str, float]|
              \item docstring: 明示されていない。
              \item このブロックが直接呼ぶ主な関数・メソッド: clip, float, to\_dict, zip
              \item 戻り値の要点: weights
              \item 主なローカル代入名: logv, raw, spec, value, weights
              \item 読み取る主なインスタンス属性: 特になし。
              \item 更新する主なインスタンス属性: 特になし。
              \item 明示的に送出する例外: 明示的raiseはない。
              \item 制御構造の規模: 条件分岐 1、ループ 1、try 0
            \end{itemize}
            \paragraph{この定義を読むためのPython構文}
            \begin{itemize}[leftmargin=1.6em]
            \item `->`以後は戻り値の型注釈であり、実行時の値を自動変換する命令ではない。
            \end{itemize}
            \paragraph{上から順の処理}
            \begin{enumerate}[leftmargin=1.8em]
            \item weights に base\_cfg.to\_dict() の結果を代入する。
\item zip(specs, vec) を順に走査し、各要素を (spec, raw) に入れて処理する。
\item   logv に float(np.clip(raw, spec.lo, spec.hi)) の結果を代入する。
\item   value に 10.0 ** logv の結果を代入する。
\item   条件 value < spec.threshold を判定し、真なら内部処理を行う。
\item     value に 0.0 の結果を代入する。
\item   weights[spec.name] に float(value) の結果を代入する。
\item weights を返す。
            \end{enumerate}
            \paragraph{代表コード断片}
            \begin{lstlisting}[language=Python]
def vector_to_weights(specs: List[TermSpec], vec: np.ndarray, base_cfg: UpperCostConfig) -> Dict[str, float]:
    weights = base_cfg.to_dict()
    for spec, raw in zip(specs, vec):
        logv = float(np.clip(raw, spec.lo, spec.hi))
        value = 10.0 ** logv
        if value < spec.threshold:
            value = 0.0
        weights[spec.name] = float(value)
    return weights
\end{lstlisting}

\subsection{L347 関数 mirror\_legacy\_weights}
            \begin{itemize}[leftmargin=1.6em]
              \item 行範囲: L347--L357
              \item 所属: module直下
              \item 定義: \path|mirror_legacy_weights(cfg: dict, upper_cost: Dict[str, float]) -> None|
              \item docstring: 明示されていない。
              \item このブロックが直接呼ぶ主な関数・メソッド: float, get, setdefault
              \item 戻り値の要点: 戻り値が無いか、return 文が明示されていない。
              \item 主なローカル代入名: mpc
              \item 読み取る主なインスタンス属性: 特になし。
              \item 更新する主なインスタンス属性: 特になし。
              \item 明示的に送出する例外: 明示的raiseはない。
              \item 制御構造の規模: 条件分岐 0、ループ 0、try 0
            \end{itemize}
            \paragraph{この定義を読むためのPython構文}
            \begin{itemize}[leftmargin=1.6em]
            \item `->`以後は戻り値の型注釈であり、実行時の値を自動変換する命令ではない。
            \end{itemize}
            \paragraph{上から順の処理}
            \begin{enumerate}[leftmargin=1.8em]
            \item mpc に cfg.setdefault('mpc', \{\}) の結果を代入する。
\item mpc['upper\_cost'] に upper\_cost の結果を代入する。
\item mpc['w\_dv'] に float(upper\_cost.get('w\_speed\_smooth', mpc.get('w\_dv', 30.0))) の結果を代入する。
\item mpc['w\_dv\_limit'] に float(upper\_cost.get('w\_dv\_limit', mpc.get('w\_dv\_limit', 2.0))) の結果を代入する。
\item mpc['w\_speed\_limit'] に float(upper\_cost.get('w\_speed\_limit', mpc.get('w\_speed\_limit', 50.0))) の結果を代入する。
\item mpc['w\_current'] に float(upper\_cost.get('w\_current\_sq', mpc.get('w\_current', 0.01))) の結果を代入する。
\item mpc['w\_T'] に float(upper\_cost.get('w\_temp', mpc.get('w\_T', 5.0))) の結果を代入する。
\item mpc['w\_soc\_day\_max'] に float(upper\_cost.get('w\_soc\_day\_max', mpc.get('w\_soc\_day\_max', 10000.0))) の結果を代入する。
\item mpc['w\_soc\_day\_track'] に float(upper\_cost.get('w\_soc\_day\_track', mpc.get('w\_soc\_day\_track', 0.0))) の結果を代入する。
\item mpc['w\_soc\_terminal'] に float(upper\_cost.get('w\_soc\_terminal', mpc.get('w\_soc\_terminal', 0.0))) の結果を代入する。
            \end{enumerate}
            \paragraph{代表コード断片}
            \begin{lstlisting}[language=Python]
def mirror_legacy_weights(cfg: dict, upper_cost: Dict[str, float]) -> None:
    mpc = cfg.setdefault("mpc", {})
    mpc["upper_cost"] = upper_cost
    mpc["w_dv"] = float(upper_cost.get("w_speed_smooth", mpc.get("w_dv", 30.0)))
    mpc["w_dv_limit"] = float(upper_cost.get("w_dv_limit", mpc.get("w_dv_limit", 2.0)))
    mpc["w_speed_limit"] = float(upper_cost.get("w_speed_limit", mpc.get("w_speed_limit", 50.0)))
    mpc["w_current"] = float(upper_cost.get("w_current_sq", mpc.get("w_current", 0.01)))
    mpc["w_T"] = float(upper_cost.get("w_temp", mpc.get("w_T", 5.0)))
    mpc["w_soc_day_max"] = float(upper_cost.get("w_soc_day_max", mpc.get("w_soc_day_max", 1.0e4)))
    mpc["w_soc_day_track"] = float(upper_cost.get("w_soc_day_track", mpc.get("w_soc_day_track", 0.0)))
    mpc["w_soc_terminal"] = float(upper_cost.get("w_soc_terminal", mpc.get("w_soc_terminal", 0.0)))
\end{lstlisting}

\subsection{L360 関数 build\_reference\_free\_profile}
            \begin{itemize}[leftmargin=1.6em]
              \item 行範囲: L360--L394
              \item 所属: module直下
              \item 定義: \path|build_reference_free_profile(profile_yaml: Path, output_dir: Path, *, disable_uncertainty_reserve: bool = False) -> tuple[Path, str]|
              \item docstring: 明示されていない。
              \item このブロックが直接呼ぶ主な関数・メソッド: Path, append, fspath, is\_absolute, isinstance, items, list, pop, read\_yaml, resolve, setdefault, str
              \item 戻り値の要点: (out\_path, removed\_reference)
              \item 主なローカル代入名: candidate, cfg, key, meta, mpc, notes, out\_path, paths, ref\_cfg, removed\_reference, upper\_cost, value
              \item 読み取る主なインスタンス属性: 特になし。
              \item 更新する主なインスタンス属性: 特になし。
              \item 明示的に送出する例外: 明示的raiseはない。
              \item 制御構造の規模: 条件分岐 7、ループ 1、try 0
            \end{itemize}
            \paragraph{この定義を読むためのPython構文}
            \begin{itemize}[leftmargin=1.6em]
            \item `->`以後は戻り値の型注釈であり、実行時の値を自動変換する命令ではない。
            \end{itemize}
            \paragraph{上から順の処理}
            \begin{enumerate}[leftmargin=1.8em]
            \item cfg に read\_yaml(profile\_yaml) の結果を代入する。
\item removed\_reference に '' の結果を代入する。
\item paths に cfg.setdefault('paths', \{\}) の結果を代入する。
\item 条件 isinstance(paths, dict) を判定し、真なら内部処理を行う。
\item   removed\_reference に str(paths.pop('progress\_reference\_csv', '') or '') の結果を代入する。
\item   list(paths.items()) を順に走査し、各要素を (key, value) に入れて処理する。
\item     条件 isinstance(value, str) and value.strip() を判定し、真なら内部処理を行う。
\item       candidate に Path(value) の結果を代入する。
\item       条件 not candidate.is\_absolute() を判定し、真なら内部処理を行う。
\item meta に cfg.setdefault('meta', \{\}) の結果を代入する。
\item notes に meta.setdefault('notes', []) の結果を代入する。
\item 条件 isinstance(notes, list) を判定し、真なら内部処理を行う。
\item   notes.append(...) を実行する。
\item mpc に cfg.setdefault('mpc', \{\}) の結果を代入する。
\item ref\_cfg に mpc.setdefault('reference\_speed\_tracking', \{\}) の結果を代入する。
\item 条件 isinstance(ref\_cfg, dict) を判定し、真なら内部処理を行う。
\item   ref\_cfg['enabled'] に False の結果を代入する。
\item upper\_cost に mpc.setdefault('upper\_cost', \{\}) の結果を代入する。
\item 条件 isinstance(upper\_cost, dict) を判定し、真なら内部処理を行う。
\item   upper\_cost['w\_progress\_lag'] に 0.0 の結果を代入する。
\item   upper\_cost['w\_progress\_terminal\_lag'] に 0.0 の結果を代入する。
\item   条件 disable\_uncertainty\_reserve を判定し、真なら内部処理を行う。
\item     upper\_cost['w\_uncertainty\_reserve'] に 0.0 の結果を代入する。
\item     upper\_cost['reserve\_soc\_per\_hour'] に 0.0 の結果を代入する。
\item     upper\_cost['reserve\_soc\_max\_extra'] に 0.0 の結果を代入する。
\item out\_path に output\_dir / 'self\_learning\_reference\_free\_profile.yaml' の結果を代入する。
\item write\_yaml(...) を実行する。
\item (out\_path, removed\_reference) を返す。
            \end{enumerate}
            \paragraph{代表コード断片}
            \begin{lstlisting}[language=Python]
def build_reference_free_profile(
    profile_yaml: Path,
    output_dir: Path,
    *,
    disable_uncertainty_reserve: bool = False,
) -> tuple[Path, str]:
    cfg = read_yaml(profile_yaml)
    removed_reference = ""
    paths = cfg.setdefault("paths", {})
    if isinstance(paths, dict):
        removed_reference = str(paths.pop("progress_reference_csv", "") or "")
        for key, value in list(paths.items()):
            if isinstance(value, str) and value.strip():
                candidate = Path(value)
                if not candidate.is_absolute():
                    paths[key] = os.fspath((profile_yaml.parent / candidate).resolve())
    meta = cfg.setdefault("meta", {})
    notes = meta.setdefault("notes", [])
    if isinstance(notes, list):
        notes.append("Self-learning tuner generated a reference-free copy for autonomous weight search.")
    mpc = cfg.setdefault("mpc", {})
    ref_cfg = mpc.setdefault("reference_speed_tracking", {})
    if isinstance(ref_cfg, dict):
        ref_cfg["enabled"] = False
    upper_cost = mpc.setdefault("upper_cost", {})
    if isinstance(upper_cost, dict):
        upper_cost["w_progress_lag"] = 0.0
        upper_cost["w_progress_terminal_lag"] = 0.0
        if disable_uncertainty_reserve:
            upper_cost["w_uncertainty_reserve"] = 0.0
            upper_cost["reserve_soc_per_hour"] = 0.0
            upper_cost["reserve_soc_max_extra"] = 0.0
    out_path = output_dir / "self_learning_reference_free_profile.yaml"
    write_yaml(out_path, cfg)
    return out_path, removed_reference
\end{lstlisting}

\subsection{L397 関数 default\_scenarios}
            \begin{itemize}[leftmargin=1.6em]
              \item 行範囲: L397--L479
              \item 所属: module直下
              \item 定義: \path|default_scenarios(profile_yaml: Path, *, mode: str = 'nominal') -> List[ScenarioSpec]|
              \item docstring: 明示されていない。
              \item このブロックが直接呼ぶ主な関数・メソッド: ScenarioSpec, float, get, isinstance, lower, max, min, read\_yaml, str, strip
              \item 戻り値の要点: [ScenarioSpec('nominal', cfg\_overrides=\{\}, cli\_overrides=\{\}, weight=0.3), ScenarioSpec('low\_solar\_high\_load', cfg\_overrides=\{'model.P\_aux': max(20.0, base\_aux + 20.0), 'model.CdA': max(base\_cda * 1.1, base\_cda + 0.005), 'model.Crr': max(base\_crr * 1.05, base\_crr + 0.0002)\}, cli\_overrides=\{'solar\_gain': 0.9, 'poa\_gain\_drive': 0.94, 'poa\_gain\_stop': 0.92\}, weight=0.2), ScenarioSpec('drag\_bias', cfg\_overrides=\{'model.P\_aux': max(10.0, base\_aux + 10.0), 'model.CdA': max(base\_cda * 1.2, base\_cda + 0.01), 'model.Crr': max(base\_crr * 1.08, base\_crr + 0.0004)\}, cli\_overrides=\{'solar\_gain': 0.97, 'poa\_gain\_drive': 0.98, 'poa\_gain\_stop': 0.98\}, weight=0.1), ScenarioSpec('low\_initial\_soc', cfg\_overrides=\{\}, cli\_overrides=\{'soc0': max(soc\_min + 0.04, base\_soc - 0.15)\}, weight=0.15), ScenarioSpec('hot\_battery', cfg\_overrides=\{\}, cli\_overrides=\{'Tb0': min(temp\_max\_c - 2.0, base\_tb\_c + 12.0)\}, weight=0.1), ScenarioSpec('driver\_lag', cfg\_overrides=\{\}, cli\_overrides=\{'exec\_tau\_sec': 4.0, 'exec\_reaction\_delay\_sec': 3.0, 'exec\_accel\_limit\_kmhps': 0.7, 'exec\_decel\_limit\_kmhps': 2.0\}, weight=0.1), ScenarioSpec('favorable\_weather', cfg\_overrides=\{'model.P\_aux': max(0.0, base\_aux - 10.0)\}, cli\_overrides=\{'solar\_gain': 1.08, 'poa\_gain\_drive': 1.04, 'poa\_gain\_stop': 1.04\}, weight=0.05)] / [ScenarioSpec('nominal', cfg\_overrides=\{\}, cli\_overrides=\{\}, weight=1.0)]
              \item 主なローカル代入名: base\_aux, base\_cda, base\_crr, base\_soc, base\_tb\_c, cfg, model, simulation, soc\_min, temp\_max\_c
              \item 読み取る主なインスタンス属性: 特になし。
              \item 更新する主なインスタンス属性: 特になし。
              \item 明示的に送出する例外: 明示的raiseはない。
              \item 制御構造の規模: 条件分岐 1、ループ 0、try 0
            \end{itemize}
            \paragraph{この定義を読むためのPython構文}
            \begin{itemize}[leftmargin=1.6em]
            \item `->`以後は戻り値の型注釈であり、実行時の値を自動変換する命令ではない。
            \end{itemize}
            \paragraph{上から順の処理}
            \begin{enumerate}[leftmargin=1.8em]
            \item 条件 str(mode).strip().lower() == 'nominal' を判定し、真なら内部処理を行う。
\item   [ScenarioSpec('nominal', cfg\_overrides=\{\}, cli\_overrides=\{\}, weight=1.0)] を返す。
\item cfg に read\_yaml(profile\_yaml) の結果を代入する。
\item model に cfg.get('model', \{\}) if isinstance(cfg, dict) else \{\} の結果を代入する。
\item base\_cda に float(model.get('CdA', 0.08)) の結果を代入する。
\item base\_crr に float(model.get('Crr', 0.008)) の結果を代入する。
\item base\_aux に float(model.get('P\_aux', 0.0)) の結果を代入する。
\item simulation に cfg.get('simulation', \{\}) if isinstance(cfg, dict) else \{\} の結果を代入する。
\item base\_soc に float(simulation.get('soc0', 0.99)) の結果を代入する。
\item base\_tb\_c に float(simulation.get('Tb0', 30.0)) の結果を代入する。
\item soc\_min に float(model.get('soc\_min', 0.05)) の結果を代入する。
\item temp\_max\_c に float(model.get('T\_max', 55.0)) の結果を代入する。
\item [ScenarioSpec('nominal', cfg\_overrides=\{\}, cli\_overrides=\{\}, weight=0.3), ScenarioSpec('low\_solar\_high\_load', cfg\_overrides=\{'model.P\_aux': max(20.0, base\_aux + 20.0), 'model.CdA': max(base\_cda * 1.1, base\_cda + 0.005), 'model.Crr': max(base\_crr * 1.05, base\_crr + 0.0002)\}, cli\_overrides=\{'solar\_gain': 0.9, 'poa\_gain\_drive': 0.94, 'poa\_gain\_stop': 0.92\}, weight=0.2), ScenarioSpec('drag\_bias', cfg\_overrides=\{'model.P\_aux': max(10.0, base\_aux + 10.0), 'model.CdA': max(base\_cda * 1.2, base\_cda + 0.01), 'model.Crr': max(base\_crr * 1.08, base\_crr + 0.0004)\}, cli\_overrides=\{'solar\_gain': 0.97, 'poa\_gain\_drive': 0.98, 'poa\_gain\_stop': 0.98\}, weight=0.1), ScenarioSpec('low\_initial\_soc', cfg\_overrides=\{\}, cli\_overrides=\{'soc0': max(soc\_min + 0.04, base\_soc - 0.15)\}, weight=0.15), ScenarioSpec('hot\_battery', cfg\_overrides=\{\}, cli\_overrides=\{'Tb0': min(temp\_max\_c - 2.0, base\_tb\_c + 12.0)\}, weight=0.1), ScenarioSpec('driver\_lag', cfg\_overrides=\{\}, cli\_overrides=\{'exec\_tau\_sec': 4.0, 'exec\_reaction\_delay\_sec': 3.0, 'exec\_accel\_limit\_kmhps': 0.7, 'exec\_decel\_limit\_kmhps': 2.0\}, weight=0.1), ScenarioSpec('favorable\_weather', cfg\_overrides=\{'model.P\_aux': max(0.0, base\_aux - 10.0)\}, cli\_overrides=\{'solar\_gain': 1.08, 'poa\_gain\_drive': 1.04, 'poa\_gain\_stop': 1.04\}, weight=0.05)] を返す。
            \end{enumerate}
            \paragraph{代表コード断片}
            \begin{lstlisting}[language=Python]
def default_scenarios(profile_yaml: Path, *, mode: str = "nominal") -> List[ScenarioSpec]:
    if str(mode).strip().lower() == "nominal":
        return [ScenarioSpec("nominal", cfg_overrides={}, cli_overrides={}, weight=1.0)]
    cfg = read_yaml(profile_yaml)
    model = cfg.get("model", {}) if isinstance(cfg, dict) else {}
    base_cda = float(model.get("CdA", 0.08))
    base_crr = float(model.get("Crr", 0.008))
    base_aux = float(model.get("P_aux", 0.0))
    simulation = cfg.get("simulation", {}) if isinstance(cfg, dict) else {}
    base_soc = float(simulation.get("soc0", 0.99))
    base_tb_c = float(simulation.get("Tb0", 30.0))
    soc_min = float(model.get("soc_min", 0.05))
    temp_max_c = float(model.get("T_max", 55.0))
    return [
        ScenarioSpec("nominal", cfg_overrides={}, cli_overrides={}, weight=0.30),
        ScenarioSpec(
            "low_solar_high_load",
            cfg_overrides={
                "model.P_aux": max(20.0, base_aux + 20.0),
                "model.CdA": max(base_cda * 1.10, base_cda + 0.005),
                "model.Crr": max(base_crr * 1.05, base_crr + 0.0002),
            },
            cli_overrides={
                "solar_gain": 0.90,
                "poa_gain_drive": 0.94,
                "poa_gain_stop": 0.92,
            },
            weight=0.20,
        ),
        ScenarioSpec(
            "drag_bias",
            cfg_overrides={
                "model.P_aux": max(10.0, base_aux + 10.0),
                "model.CdA": max(base_cda * 1.20, base_cda + 0.010),
                "model.Crr": max(base_crr * 1.08, base_crr + 0.0004),
...
\end{lstlisting}

\subsection{L482 関数 run\_single\_scenario}
            \begin{itemize}[leftmargin=1.6em]
              \item 行範囲: L482--L610
              \item 所属: module直下
              \item 定義: \path|run_single_scenario(profile_yaml: Path, output_dir: Path, candidate_name: str, scenario: ScenarioSpec, upper_cost: Dict[str, float], cfg_overrides: Dict[str, object], cli_overrides: Dict[str, object]) -> Dict[str, object]|
              \item docstring: 明示されていない。
              \item このブロックが直接呼ぶ主な関数・メソッド: DataFrame, Path, dict, dumps, ensure\_dir, evaluate\_simulation, exists, extend, failed\_scenario\_result, float, format\_override\_value, fspath
              \item 戻り値の要点: result
              \item 主なローカル代入名: cmd, detail\_csv, detail\_df, key, latest\_manifest\_json, log\_f, merged\_cfg, merged\_cli, out\_csv, plan\_csv, proc, report\_html, resolved\_yaml, result, scenario\_dir, sim\_df, sim\_log\_path, summary, summary\_json, value
              \item 読み取る主なインスタンス属性: 特になし。
              \item 更新する主なインスタンス属性: 特になし。
              \item 明示的に送出する例外: 明示的raiseはない。
              \item 制御構造の規模: 条件分岐 1、ループ 3、try 1
            \end{itemize}
            \paragraph{この定義を読むためのPython構文}
            \begin{itemize}[leftmargin=1.6em]
            \item `->`以後は戻り値の型注釈であり、実行時の値を自動変換する命令ではない。
\item with文はcontext managerに開始・終了処理を任せ、ファイルなどの資源を確実に解放する。
\item try文は例外が起き得る処理と、異常時または終了時の経路を分ける。
            \end{itemize}
            \paragraph{上から順の処理}
            \begin{enumerate}[leftmargin=1.8em]
            \item scenario\_dir に output\_dir / 'candidates' / candidate\_name / scenario.name の結果を代入する。
\item ensure\_dir(...) を実行する。
\item out\_csv に scenario\_dir / 'simulation.csv' の結果を代入する。
\item detail\_csv に scenario\_dir / 'simulation\_detail.csv' の結果を代入する。
\item plan\_csv に scenario\_dir / 'upper\_plan.csv' の結果を代入する。
\item report\_html に scenario\_dir / 'simulation\_report.html' の結果を代入する。
\item summary\_json に scenario\_dir / 'summary.json' の結果を代入する。
\item resolved\_yaml に scenario\_dir / 'resolved.yaml' の結果を代入する。
\item latest\_manifest\_json に scenario\_dir / 'latest\_manifest.json' の結果を代入する。
\item sim\_log\_path に scenario\_dir / 'solar\_sim\_console.log' の結果を代入する。
\item cmd に [os.fspath(Path(sys.executable)), os.fspath(ROOT / 'scripts' / 'solar\_sim.py'), '--profile\_yaml', os.fspath(profile\_yaml), '--out\_csv', os.fspath(out\_csv), '--out\_detail\_csv', os.fspath(detail\_csv), '--out\_plan\_csv', os.fspath(plan\_csv), '--report\_html', os.fspath(report\_html), '--summary\_json', os.fspath(summary\_json), '--resolved\_yaml', os.fspath(resolved\_yaml), '--latest\_manifest\_json', os.fspath(latest\_manifest\_json)] の結果を代入する。
\item merged\_cli に dict(cli\_overrides) の結果を代入する。
\item merged\_cli.update(...) を実行する。
\item merged\_cli.items() を順に走査し、各要素を (key, value) に入れて処理する。
\item   cmd.extend(...) を実行する。
\item merged\_cfg に dict(cfg\_overrides) の結果を代入する。
\item merged\_cfg.update(...) を実行する。
\item upper\_cost.items() を順に走査し、各要素を (key, value) に入れて処理する。
\item   merged\_cfg[f'mpc.upper\_cost.\{key\}'] に value の結果を代入する。
\item merged\_cfg['mpc.w\_dv'] に upper\_cost.get('w\_speed\_smooth', 0.0) の結果を代入する。
\item merged\_cfg['mpc.w\_dv\_limit'] に upper\_cost.get('w\_dv\_limit', 0.0) の結果を代入する。
\item merged\_cfg['mpc.w\_speed\_limit'] に upper\_cost.get('w\_speed\_limit', 0.0) の結果を代入する。
\item merged\_cfg['mpc.w\_current'] に upper\_cost.get('w\_current\_sq', 0.0) の結果を代入する。
\item merged\_cfg['mpc.w\_T'] に upper\_cost.get('w\_temp', 0.0) の結果を代入する。
\item merged\_cfg['mpc.w\_soc\_day\_max'] に upper\_cost.get('w\_soc\_day\_max', 0.0) の結果を代入する。
\item merged\_cfg['mpc.w\_soc\_day\_track'] に upper\_cost.get('w\_soc\_day\_track', 0.0) の結果を代入する。
\item merged\_cfg['mpc.w\_soc\_terminal'] に upper\_cost.get('w\_soc\_terminal', 0.0) の結果を代入する。
\item merged\_cfg.items() を順に走査し、各要素を (key, value) に入れて処理する。
\item   cmd.extend(...) を実行する。
\item with 文で sim\_log\_path.open('w', encoding='utf-8', newline='\textbackslash{}n') を管理しながら処理する。
\item   proc に subprocess.run(cmd, cwd=ROOT, check=False, stdout=log\_f, stderr=subprocess.STDOUT, text=True) の結果を代入する。
\item 条件 proc.returncode != 0 を判定し、真なら内部処理を行う。
\item   result に failed\_scenario\_result(scenario.name, scenario.weight, upper\_cost, error=f'solar\_sim failed with exit code \{proc.returncode\}', cfg\_overrides=merged\_cfg, cli\_overrides=merged\_cli, summary\_json=summary\_json, out\_csv=out\_csv, detail\_csv=detail\_csv, plan\_csv=plan\_csv, report\_html=report\_html, resolved\_yaml=resolved\_yaml, sim\_log\_path=sim\_log\_path) の結果を代入する。
\item   上の条件が偽の場合:
\item   例外処理を伴う try ブロックを実行する。
\item     summary に json.loads(summary\_json.read\_text(encoding='utf-8')) の結果を代入する。
\item     sim\_df に pd.read\_csv(out\_csv) の結果を代入する。
\item     detail\_df に pd.read\_csv(detail\_csv) if detail\_csv.exists() else pd.DataFrame() の結果を代入する。
\item     result に evaluate\_simulation(profile\_yaml, summary, sim\_df, detail\_df, upper\_cost) の結果を代入する。
\item     Exceptionを捕捉した場合:
\item     result に failed\_scenario\_result(scenario.name, scenario.weight, upper\_cost, error=f'postprocess failed: \{exc\}', cfg\_overrides=merged\_cfg, cli\_overrides=merged\_cli, summary\_json=summary\_json, out\_csv=out\_csv, detail\_csv=detail\_csv, plan\_csv=plan\_csv, report\_html=report\_html, resolved\_yaml=resolved\_yaml, sim\_log\_path=sim\_log\_path) の結果を代入する。
\item result.update(...) を実行する。
\item (scenario\_dir / 'eval\_metrics.json').write\_text(...) を実行する。
\item result を返す。
            \end{enumerate}
            \paragraph{代表コード断片}
            \begin{lstlisting}[language=Python]
def run_single_scenario(
    profile_yaml: Path,
    output_dir: Path,
    candidate_name: str,
    scenario: ScenarioSpec,
    upper_cost: Dict[str, float],
    cfg_overrides: Dict[str, object],
    cli_overrides: Dict[str, object],
) -> Dict[str, object]:
    scenario_dir = output_dir / "candidates" / candidate_name / scenario.name
    ensure_dir(scenario_dir)

    out_csv = scenario_dir / "simulation.csv"
    detail_csv = scenario_dir / "simulation_detail.csv"
    plan_csv = scenario_dir / "upper_plan.csv"
    report_html = scenario_dir / "simulation_report.html"
    summary_json = scenario_dir / "summary.json"
    resolved_yaml = scenario_dir / "resolved.yaml"
    latest_manifest_json = scenario_dir / "latest_manifest.json"
    sim_log_path = scenario_dir / "solar_sim_console.log"

    cmd = [
        os.fspath(Path(sys.executable)),
        os.fspath(ROOT / "scripts" / "solar_sim.py"),
        "--profile_yaml",
        os.fspath(profile_yaml),
        "--out_csv",
        os.fspath(out_csv),
        "--out_detail_csv",
        os.fspath(detail_csv),
        "--out_plan_csv",
        os.fspath(plan_csv),
        "--report_html",
        os.fspath(report_html),
        "--summary_json",
...
\end{lstlisting}

\subsection{L613 関数 evaluate\_simulation}
            \begin{itemize}[leftmargin=1.6em]
              \item 行範囲: L613--L756
              \item 所属: module直下
              \item 定義: \path|evaluate_simulation(profile_yaml: Path, summary: Dict[str, object], sim_df: pd.DataFrame, detail_df: pd.DataFrame, upper_cost: Dict[str, float]) -> Dict[str, object]|
              \item docstring: 明示されていない。
              \item このブロックが直接呼ぶ主な関数・メソッド: Series, UpperCostConfig, abs, active\_upper\_cost\_terms, astype, bool, diff, exists, fillna, float, from\_yaml, fspath
              \item 戻り値の要点: \{'score': float(score), 'final\_distance\_km': float(summary['final\_distance\_km']), 'avg\_speed\_kmh': float(summary.get('avg\_speed\_kmh', 0.0)), 'min\_soc': min\_soc, 'final\_soc': final\_soc, 'elapsed\_hours': elapsed\_hours, 'cpu\_sec': float(summary.get('cpu\_sec', 0.0)), 'finish\_reached': finish\_reached, 'scenario\_feasible': not infeasibility\_reasons, 'scenario\_feasibility\_checks': feasibility\_checks, 'scenario\_infeasibility\_reasons': infeasibility\_reasons, 'oscillation\_mean\_abs\_dv\_kmh': osc, 'oscillation\_p95\_abs\_dv\_kmh': osc\_p95, 'current\_rms\_a': current\_rms\_a, 'pack\_slew\_rms\_kw': pack\_slew\_rms\_kw, 'high\_speed\_h': high\_speed\_h, 'daylight\_stop\_h': daylight\_stop\_h, 'daylight\_full\_soc\_h': daylight\_full\_soc\_h, 'unused\_finish\_soc': unused\_finish\_soc, 'finish\_soc\_target': finish\_soc\_target, 'terminal\_soc\_error': terminal\_soc\_error, 'active\_term\_count': int(len(active\_terms)), 'active\_terms': active\_terms\}
              \item 主なローカル代入名: active\_pair, active\_terms, current\_rms\_a, current\_vals, daylight\_full\_soc\_h, daylight\_mask, daylight\_stop\_h, drive\_mask, dt\_hours, elapsed\_hours, feasibility\_checks, final\_soc, finish\_bonus, finish\_reached, finish\_soc\_target, high\_speed\_h, high\_speed\_mask, infeasibility\_reasons, min\_soc, mpc\_cfg
              \item 読み取る主なインスタンス属性: 特になし。
              \item 更新する主なインスタンス属性: 特になし。
              \item 明示的に送出する例外: 明示的raiseはない。
              \item 制御構造の規模: 条件分岐 9、ループ 0、try 0
            \end{itemize}
            \paragraph{この定義を読むためのPython構文}
            \begin{itemize}[leftmargin=1.6em]
            \item `->`以後は戻り値の型注釈であり、実行時の値を自動変換する命令ではない。
\item lambdaは名前を付けずに短い関数オブジェクトを作る構文である。
\item 内包表記は、反復・条件・値生成を一つの式で記述してcollectionを作る。
            \end{itemize}
            \paragraph{上から順の処理}
            \begin{enumerate}[leftmargin=1.8em]
            \item speed\_vals に speed\_series(sim\_df).to\_numpy(dtype=float) の結果を代入する。
\item 条件 'time\_utc' in sim\_df.columns and len(sim\_df) >= 2 を判定し、真なら内部処理を行う。
\item   t\_series に pd.to\_datetime(sim\_df['time\_utc'], format='mixed', utc=True, errors='coerce') の結果を代入する。
\item   dt\_hours に t\_series.diff().dt.total\_seconds().fillna(t\_series.diff().dt.total\_seconds().median()).fillna(0.0) / 3600.0 の結果を代入する。
\item   上の条件が偽の場合:
\item   t\_series に pd.Series([pd.NaT] * len(sim\_df)) の結果を代入する。
\item   dt\_hours に pd.Series(np.zeros(len(sim\_df), dtype=float)) の結果を代入する。
\item profile\_cfg に read\_yaml(profile\_yaml) の結果を代入する。
\item schedule に None の結果を代入する。
\item schedule\_rel に ((profile\_cfg.get('paths', \{\}) if isinstance(profile\_cfg, dict) else \{\}) or \{\}).get('drive\_schedule\_yaml') の結果を代入する。
\item 条件 schedule\_rel を判定し、真なら内部処理を行う。
\item   schedule\_path に (profile\_yaml.parent / schedule\_rel).resolve() の結果を代入する。
\item   条件 schedule\_path.exists() を判定し、真なら内部処理を行う。
\item     schedule に DriveSchedule.from\_yaml(os.fspath(schedule\_path)) の結果を代入する。
\item daylight\_mask に sim\_df.get('G\_poa', pd.Series(np.zeros(len(sim\_df), dtype=float))).astype(float) >= 250.0 の結果を代入する。
\item stopped\_mask に speed\_series(sim\_df) <= 1.0 の結果を代入する。
\item soc\_mask に sim\_df.get('soc', pd.Series(np.zeros(len(sim\_df), dtype=float))).astype(float) >= 0.95 の結果を代入する。
\item 条件 schedule is not None and len(t\_series) == len(sim\_df) を判定し、真なら内部処理を行う。
\item   drive\_mask に t\_series.map(lambda ts: bool(pd.notna(ts) and schedule.is\_drive\_time(ts.to\_pydatetime()))) の結果を代入する。
\item   上の条件が偽の場合:
\item   drive\_mask に pd.Series(np.ones(len(sim\_df), dtype=bool)) の結果を代入する。
\item 条件 len(speed\_vals) >= 2 を判定し、真なら内部処理を行う。
\item   active\_pair に drive\_mask.to\_numpy(dtype=bool)[1:] \& drive\_mask.to\_numpy(dtype=bool)[:-1] \& (speed\_vals[1:] > 1.0) \& (speed\_vals[:-1] > 1.0) の結果を代入する。
\item   speed\_steps に np.abs(np.diff(speed\_vals))[active\_pair] の結果を代入する。
\item   上の条件が偽の場合:
\item   speed\_steps に np.zeros(0, dtype=float) の結果を代入する。
\item osc に float(np.mean(speed\_steps)) if len(speed\_steps) else 0.0 の結果を代入する。
\item osc\_p95 に float(np.percentile(speed\_steps, 95.0)) if len(speed\_steps) else 0.0 の結果を代入する。
\item daylight\_stop\_h に float(dt\_hours[drive\_mask \& daylight\_mask \& stopped\_mask].sum()) if len(dt\_hours) == len(sim\_df) else 0.0 の結果を代入する。
\item daylight\_full\_soc\_h に float(dt\_hours[drive\_mask \& daylight\_mask \& soc\_mask].sum()) if len(dt\_hours) == len(sim\_df) else 0.0 の結果を代入する。
\item high\_speed\_mask に speed\_series(sim\_df) >= 85.0 の結果を代入する。
\item high\_speed\_h に float(dt\_hours[drive\_mask \& high\_speed\_mask].sum()) if len(dt\_hours) == len(sim\_df) else 0.0 の結果を代入する。
\item 条件 not detail\_df.empty and 'I' in detail\_df.columns を判定し、真なら内部処理を行う。
\item   current\_vals に detail\_df['I'].to\_numpy(dtype=float) の結果を代入する。
\item   current\_rms\_a に float(np.sqrt(np.mean(np.square(current\_vals)))) if len(current\_vals) else 0.0 の結果を代入する。
\item   上の条件が偽の場合:
\item   current\_rms\_a に 0.0 の結果を代入する。
\item 条件 not detail\_df.empty and 'P\_pack' in detail\_df.columns and (len(detail\_df) >= 2) を判定し、真なら内部処理を行う。
\item   pack\_slew\_kw に np.diff(detail\_df['P\_pack'].to\_numpy(dtype=float)) / 1000.0 の結果を代入する。
\item   pack\_slew\_rms\_kw に float(np.sqrt(np.mean(np.square(pack\_slew\_kw)))) if len(pack\_slew\_kw) else 0.0 の結果を代入する。
\item   上の条件が偽の場合:
\item   pack\_slew\_rms\_kw に 0.0 の結果を代入する。
\item final\_soc に float(summary.get('final\_soc', 0.0)) の結果を代入する。
\item mpc\_cfg に profile\_cfg.get('mpc', \{\}) if isinstance(profile\_cfg, dict) else \{\} の結果を代入する。
\item finish\_soc\_target に float(mpc\_cfg.get('soc\_finish\_target', 0.12) or 0.12) の結果を代入する。
\item 条件 finish\_soc\_target <= 0.0 を判定し、真なら内部処理を行う。
\item   finish\_soc\_target に 0.12 の結果を代入する。
\item unused\_finish\_soc に max(0.0, final\_soc - finish\_soc\_target) の結果を代入する。
\item terminal\_soc\_error に abs(final\_soc - finish\_soc\_target) の結果を代入する。
\item finish\_reached に bool(summary.get('finish\_reached', False)) の結果を代入する。
\item min\_soc に float(summary.get('min\_soc', 0.0)) の結果を代入する。
\item elapsed\_hours に float(summary.get('elapsed\_hours', 0.0)) の結果を代入する。
\item active\_terms に active\_upper\_cost\_terms(UpperCostConfig(**upper\_cost), threshold=0.0001) の結果を代入する。
\item finish\_bonus に 10000.0 if finish\_reached else 0.0 の結果を代入する。
\item score に float(summary['final\_distance\_km']) + finish\_bonus - 2.0 * osc - 1.0 * osc\_p95 - 0.2 * current\_rms\_a - 12.0 * pack\_slew\_rms\_kw - 10.0 * high\_speed\_h - 40.0 * daylight\_stop\_h - 100.0 * daylight\_full\_soc\_h - 1000.0 * terminal\_soc\_error - 0.5 * len(active\_terms) - (30.0 * elapsed\_hours if finish\_reached else 0.0) の結果を代入する。
\item upper\_solve\_count に int(summary.get('upper\_solve\_count', 0) or 0) の結果を代入する。
\item terminal\_soc\_min に float(summary.get('terminal\_soc\_min', mpc\_cfg.get('terminal\_soc\_min', 0.0))) の結果を代入する。
\item feasibility\_checks に \{'finish\_reached': finish\_reached, 'model\_validation\_gate\_pass': bool(summary.get('model\_validation\_gate\_pass', False)), 'terminal\_soc\_target\_met': bool(summary.get('terminal\_soc\_target\_met', True)), 'minimum\_soc\_respected': bool(math.isfinite(min\_soc) and min\_soc >= terminal\_soc\_min - 1e-06), 'finite\_score': bool(math.isfinite(score))\} の結果を代入する。
\item 条件 upper\_solve\_count > 0 を判定し、真なら内部処理を行う。
\item   feasibility\_checks.update(...) を実行する。
\item infeasibility\_reasons に [name for name, passed in feasibility\_checks.items() if not bool(passed)] の結果を代入する。
\item \{'score': float(score), 'final\_distance\_km': float(summary['final\_distance\_km']), 'avg\_speed\_kmh': float(summary.get('avg\_speed\_kmh', 0.0)), 'min\_soc': min\_soc, 'final\_soc': final\_soc, 'elapsed\_hours': elapsed\_hours, 'cpu\_sec': float(summary.get('cpu\_sec', 0.0)), 'finish\_reached': finish\_reached, 'scenario\_feasible': not infeasibility\_reasons, 'scenario\_feasibility\_checks': feasibility\_checks, 'scenario\_infeasibility\_reasons': infeasibility\_reasons, 'oscillation\_mean\_abs\_dv\_kmh': osc, 'oscillation\_p95\_abs\_dv\_kmh': osc\_p95, 'current\_rms\_a': current\_rms\_a, 'pack\_slew\_rms\_kw': pack\_slew\_rms\_kw, 'high\_speed\_h': high\_speed\_h, 'daylight\_stop\_h': daylight\_stop\_h, 'daylight\_full\_soc\_h': daylight\_full\_soc\_h, 'unused\_finish\_soc': unused\_finish\_soc, 'finish\_soc\_target': finish\_soc\_target, 'terminal\_soc\_error': terminal\_soc\_error, 'active\_term\_count': int(len(active\_terms)), 'active\_terms': active\_terms\} を返す。
            \end{enumerate}
            \paragraph{代表コード断片}
            \begin{lstlisting}[language=Python]
def evaluate_simulation(
    profile_yaml: Path,
    summary: Dict[str, object],
    sim_df: pd.DataFrame,
    detail_df: pd.DataFrame,
    upper_cost: Dict[str, float],
) -> Dict[str, object]:
    speed_vals = speed_series(sim_df).to_numpy(dtype=float)

    if "time_utc" in sim_df.columns and len(sim_df) >= 2:
        t_series = pd.to_datetime(sim_df["time_utc"], format="mixed", utc=True, errors="coerce")
        dt_hours = t_series.diff().dt.total_seconds().fillna(t_series.diff().dt.total_seconds().median()).fillna(0.0) / 3600.0
    else:
        t_series = pd.Series([pd.NaT] * len(sim_df))
        dt_hours = pd.Series(np.zeros(len(sim_df), dtype=float))

    profile_cfg = read_yaml(profile_yaml)
    schedule = None
    schedule_rel = ((profile_cfg.get("paths", {}) if isinstance(profile_cfg, dict) else {}) or {}).get("drive_schedule_yaml")
    if schedule_rel:
        schedule_path = (profile_yaml.parent / schedule_rel).resolve()
        if schedule_path.exists():
            schedule = DriveSchedule.from_yaml(os.fspath(schedule_path))

    daylight_mask = sim_df.get("G_poa", pd.Series(np.zeros(len(sim_df), dtype=float))).astype(float) >= 250.0
    stopped_mask = speed_series(sim_df) <= 1.0
    soc_mask = sim_df.get("soc", pd.Series(np.zeros(len(sim_df), dtype=float))).astype(float) >= 0.95
    if schedule is not None and len(t_series) == len(sim_df):
        drive_mask = t_series.map(lambda ts: bool(pd.notna(ts) and schedule.is_drive_time(ts.to_pydatetime())))
    else:
        drive_mask = pd.Series(np.ones(len(sim_df), dtype=bool))
    # A scheduled stop/start is not controller chatter.  Score speed variation
    # only between consecutive samples that are both actively driving.
    if len(speed_vals) >= 2:
        active_pair = (
...
\end{lstlisting}

\subsection{L759 関数 aggregate\_candidate}
            \begin{itemize}[leftmargin=1.6em]
              \item 行範囲: L759--L817
              \item 所属: module直下
              \item 定義: \path!aggregate_candidate(candidate: str, scenario_results: List[Dict[str, object]], weights: Dict[str, float], risk_config: ScenarioRiskConfig | None = None) -> Dict[str, object]!
              \item docstring: 明示されていない。
              \item このブロックが直接呼ぶ主な関数・メソッド: ValueError, aggregate\_scenario\_scores, all, array, asarray, bool, dot, float, get, int, isfinite, len
              \item 戻り値の要点: result
              \item 主なローカル代入名: nominal, result, risk, row, scenario\_feasible, scenario\_scores, scenario\_weights
              \item 読み取る主なインスタンス属性: 特になし。
              \item 更新する主なインスタンス属性: 特になし。
              \item 明示的に送出する例外: ValueError('scenario\_results is empty')
              \item 制御構造の規模: 条件分岐 2、ループ 0、try 0
            \end{itemize}
            \paragraph{この定義を読むためのPython構文}
            \begin{itemize}[leftmargin=1.6em]
            \item `->`以後は戻り値の型注釈であり、実行時の値を自動変換する命令ではない。
\item 内包表記は、反復・条件・値生成を一つの式で記述してcollectionを作る。
            \end{itemize}
            \paragraph{上から順の処理}
            \begin{enumerate}[leftmargin=1.8em]
            \item 条件 not scenario\_results を判定し、真なら内部処理を行う。
\item   ValueError('scenario\_results is empty') を送出する。
\item scenario\_weights に np.array([max(0.0, float(row.get('scenario\_weight', 0.0))) for row in scenario\_results], dtype=float) の結果を代入する。
\item 条件 not np.isfinite(scenario\_weights).all() or scenario\_weights.sum() <= 0.0 を判定し、真なら内部処理を行う。
\item   scenario\_weights に np.ones(len(scenario\_results), dtype=float) の結果を代入する。
\item scenario\_scores に np.array([float(row['score']) for row in scenario\_results], dtype=float) の結果を代入する。
\item scenario\_feasible に np.array([bool(row.get('scenario\_feasible', False)) for row in scenario\_results], dtype=bool) の結果を代入する。
\item risk に aggregate\_scenario\_scores(scenario\_scores, scenario\_weights, scenario\_feasible, config=risk\_config) の結果を代入する。
\item scenario\_weights に np.asarray(risk['normalized\_scenario\_weights'], dtype=float) の結果を代入する。
\item nominal に next((row for row in scenario\_results if row.get('scenario') == 'nominal'), scenario\_results[0]) の結果を代入する。
\item result に \{'candidate': candidate, **risk, 'final\_distance\_km': float(np.dot(scenario\_weights, np.array([float(row['final\_distance\_km']) for row in scenario\_results]))), 'final\_distance\_worst\_km': float(min((float(row['final\_distance\_km']) for row in scenario\_results))), 'avg\_speed\_kmh': float(np.dot(scenario\_weights, np.array([float(row['avg\_speed\_kmh']) for row in scenario\_results]))), 'min\_soc': float(min((float(row['min\_soc']) for row in scenario\_results))), 'final\_soc': float(np.dot(scenario\_weights, np.array([float(row['final\_soc']) for row in scenario\_results]))), 'finish\_reached': bool(all((bool(row['finish\_reached']) for row in scenario\_results))), 'model\_validation\_gate\_pass\_all': bool(all((bool(row.get('model\_validation\_gate\_pass', False)) for row in scenario\_results))), 'oscillation\_mean\_abs\_dv\_kmh': float(np.dot(scenario\_weights, np.array([float(row['oscillation\_mean\_abs\_dv\_kmh']) for row in scenario\_results]))), 'oscillation\_p95\_abs\_dv\_kmh': float(np.dot(scenario\_weights, np.array([float(row['oscillation\_p95\_abs\_dv\_kmh']) for row in scenario\_results]))), 'current\_rms\_a': float(np.dot(scenario\_weights, np.array([float(row['current\_rms\_a']) for row in scenario\_results]))), 'pack\_slew\_rms\_kw': float(np.dot(scenario\_weights, np.array([float(row['pack\_slew\_rms\_kw']) for row in scenario\_results]))), 'high\_speed\_h': float(np.dot(scenario\_weights, np.array([float(row['high\_speed\_h']) for row in scenario\_results]))), 'daylight\_stop\_h': float(np.dot(scenario\_weights, np.array([float(row['daylight\_stop\_h']) for row in scenario\_results]))), 'daylight\_full\_soc\_h': float(np.dot(scenario\_weights, np.array([float(row['daylight\_full\_soc\_h']) for row in scenario\_results]))), 'unused\_finish\_soc': float(np.dot(scenario\_weights, np.array([float(row['unused\_finish\_soc']) for row in scenario\_results]))), 'finish\_soc\_target': float(np.dot(scenario\_weights, np.array([float(row['finish\_soc\_target']) for row in scenario\_results]))), 'terminal\_soc\_error': float(np.dot(scenario\_weights, np.array([float(row['terminal\_soc\_error']) for row in scenario\_results]))), 'elapsed\_hours': float(np.dot(scenario\_weights, np.array([float(row['elapsed\_hours']) for row in scenario\_results]))), 'cpu\_sec': float(sum((float(row['cpu\_sec']) for row in scenario\_results))), 'active\_term\_count': int(round(np.dot(scenario\_weights, np.array([float(row['active\_term\_count']) for row in scenario\_results])))), 'weights': weights, 'scenario\_results': scenario\_results, 'nominal\_out\_csv': nominal.get('out\_csv', ''), 'nominal\_detail\_csv': nominal.get('detail\_csv', '')\} の結果を代入する。
\item result を返す。
            \end{enumerate}
            \paragraph{代表コード断片}
            \begin{lstlisting}[language=Python]
def aggregate_candidate(
    candidate: str,
    scenario_results: List[Dict[str, object]],
    weights: Dict[str, float],
    risk_config: ScenarioRiskConfig | None = None,
) -> Dict[str, object]:
    if not scenario_results:
        raise ValueError("scenario_results is empty")
    scenario_weights = np.array([max(0.0, float(row.get("scenario_weight", 0.0))) for row in scenario_results], dtype=float)
    if not np.isfinite(scenario_weights).all() or scenario_weights.sum() <= 0.0:
        scenario_weights = np.ones(len(scenario_results), dtype=float)
    scenario_scores = np.array([float(row["score"]) for row in scenario_results], dtype=float)
    scenario_feasible = np.array(
        [bool(row.get("scenario_feasible", False)) for row in scenario_results],
        dtype=bool,
    )
    risk = aggregate_scenario_scores(
        scenario_scores,
        scenario_weights,
        scenario_feasible,
        config=risk_config,
    )
    scenario_weights = np.asarray(risk["normalized_scenario_weights"], dtype=float)

    nominal = next((row for row in scenario_results if row.get("scenario") == "nominal"), scenario_results[0])
    result = {
        "candidate": candidate,
        **risk,
        "final_distance_km": float(np.dot(scenario_weights, np.array([float(row["final_distance_km"]) for row in scenario_results]))),
        "final_distance_worst_km": float(min(float(row["final_distance_km"]) for row in scenario_results)),
        "avg_speed_kmh": float(np.dot(scenario_weights, np.array([float(row["avg_speed_kmh"]) for row in scenario_results]))),
        "min_soc": float(min(float(row["min_soc"]) for row in scenario_results)),
        "final_soc": float(np.dot(scenario_weights, np.array([float(row["final_soc"]) for row in scenario_results]))),
        "finish_reached": bool(all(bool(row["finish_reached"]) for row in scenario_results)),
        "model_validation_gate_pass_all": bool(
...
\end{lstlisting}

\subsection{L820 関数 run\_candidate}
            \begin{itemize}[leftmargin=1.6em]
              \item 行範囲: L820--L855
              \item 所属: module直下
              \item 定義: \path!run_candidate(profile_yaml: Path, output_dir: Path, candidate_name: str, upper_cost: Dict[str, float], cfg_overrides: Dict[str, object], cli_overrides: Dict[str, object], scenarios: List[ScenarioSpec], risk_config: ScenarioRiskConfig | None = None) -> Dict[str, object]!
              \item docstring: 明示されていない。
              \item このブロックが直接呼ぶ主な関数・メソッド: aggregate\_candidate, canonical\_runtime\_weights, dumps, ensure\_dir, run\_single\_scenario, write\_text
              \item 戻り値の要点: result
              \item 主なローカル代入名: cand\_dir, result, runtime\_upper\_cost, scenario, scenario\_results
              \item 読み取る主なインスタンス属性: 特になし。
              \item 更新する主なインスタンス属性: 特になし。
              \item 明示的に送出する例外: 明示的raiseはない。
              \item 制御構造の規模: 条件分岐 0、ループ 0、try 0
            \end{itemize}
            \paragraph{この定義を読むためのPython構文}
            \begin{itemize}[leftmargin=1.6em]
            \item `->`以後は戻り値の型注釈であり、実行時の値を自動変換する命令ではない。
\item 内包表記は、反復・条件・値生成を一つの式で記述してcollectionを作る。
            \end{itemize}
            \paragraph{上から順の処理}
            \begin{enumerate}[leftmargin=1.8em]
            \item runtime\_upper\_cost に canonical\_runtime\_weights(upper\_cost) の結果を代入する。
\item scenario\_results に [run\_single\_scenario(profile\_yaml, output\_dir, candidate\_name, scenario, runtime\_upper\_cost, cfg\_overrides, cli\_overrides) for scenario in scenarios] の結果を代入する。
\item result に aggregate\_candidate(candidate\_name, scenario\_results, runtime\_upper\_cost, risk\_config=risk\_config) の結果を代入する。
\item cand\_dir に output\_dir / 'candidates' / candidate\_name の結果を代入する。
\item ensure\_dir(...) を実行する。
\item (cand\_dir / 'aggregate\_metrics.json').write\_text(...) を実行する。
\item result を返す。
            \end{enumerate}
            \paragraph{代表コード断片}
            \begin{lstlisting}[language=Python]
def run_candidate(
    profile_yaml: Path,
    output_dir: Path,
    candidate_name: str,
    upper_cost: Dict[str, float],
    cfg_overrides: Dict[str, object],
    cli_overrides: Dict[str, object],
    scenarios: List[ScenarioSpec],
    risk_config: ScenarioRiskConfig | None = None,
) -> Dict[str, object]:
    runtime_upper_cost = canonical_runtime_weights(upper_cost)
    scenario_results = [
        run_single_scenario(
            profile_yaml,
            output_dir,
            candidate_name,
            scenario,
            runtime_upper_cost,
            cfg_overrides,
            cli_overrides,
        )
        for scenario in scenarios
    ]
    result = aggregate_candidate(
        candidate_name,
        scenario_results,
        runtime_upper_cost,
        risk_config=risk_config,
    )
    cand_dir = output_dir / "candidates" / candidate_name
    ensure_dir(cand_dir)
    (cand_dir / "aggregate_metrics.json").write_text(
        json.dumps(result, indent=2, ensure_ascii=False),
        encoding="utf-8",
    )
...
\end{lstlisting}

\subsection{L858 関数 save\_trial\_checkpoint}
            \begin{itemize}[leftmargin=1.6em]
              \item 行範囲: L858--L865
              \item 所属: module直下
              \item 定義: \path!save_trial_checkpoint(output_dir: Path, trials: List[Dict[str, object]], best_result: Dict[str, object] | None) -> None!
              \item docstring: 明示されていない。
              \item このブロックが直接呼ぶ主な関数・メソッド: DataFrame, float, get, to\_csv, write\_yaml
              \item 戻り値の要点: 戻り値が無いか、return 文が明示されていない。
              \item 主なローカル代入名: 特になし。
              \item 読み取る主なインスタンス属性: 特になし。
              \item 更新する主なインスタンス属性: 特になし。
              \item 明示的に送出する例外: 明示的raiseはない。
              \item 制御構造の規模: 条件分岐 2、ループ 0、try 0
            \end{itemize}
            \paragraph{この定義を読むためのPython構文}
            \begin{itemize}[leftmargin=1.6em]
            \item `->`以後は戻り値の型注釈であり、実行時の値を自動変換する命令ではない。
            \end{itemize}
            \paragraph{上から順の処理}
            \begin{enumerate}[leftmargin=1.8em]
            \item 条件 trials を判定し、真なら内部処理を行う。
\item   pd.DataFrame(trials).to\_csv(...) を実行する。
\item 条件 best\_result is not None を判定し、真なら内部処理を行う。
\item   write\_yaml(...) を実行する。
            \end{enumerate}
            \paragraph{代表コード断片}
            \begin{lstlisting}[language=Python]
def save_trial_checkpoint(output_dir: Path, trials: List[Dict[str, object]], best_result: Dict[str, object] | None) -> None:
    if trials:
        pd.DataFrame(trials).to_csv(output_dir / "trial_results_partial.csv", index=False)
    if best_result is not None:
        write_yaml(
            output_dir / "best_upper_cost_partial.yaml",
            {"upper_cost": best_result.get("weights", {}), "score": float(best_result.get("score", 0.0))},
        )
\end{lstlisting}

\subsection{L868 関数 csv\_row\_count}
            \begin{itemize}[leftmargin=1.6em]
              \item 行範囲: L868--L870
              \item 所属: module直下
              \item 定義: \path|csv_row_count(path: Path) -> int|
              \item docstring: 明示されていない。
              \item このブロックが直接呼ぶ主な関数・メソッド: max, open, sum
              \item 戻り値の要点: max(0, sum((1 for \_ in f)) - 1)
              \item 主なローカル代入名: \_, f
              \item 読み取る主なインスタンス属性: 特になし。
              \item 更新する主なインスタンス属性: 特になし。
              \item 明示的に送出する例外: 明示的raiseはない。
              \item 制御構造の規模: 条件分岐 0、ループ 0、try 0
            \end{itemize}
            \paragraph{この定義を読むためのPython構文}
            \begin{itemize}[leftmargin=1.6em]
            \item `->`以後は戻り値の型注釈であり、実行時の値を自動変換する命令ではない。
\item 内包表記は、反復・条件・値生成を一つの式で記述してcollectionを作る。
\item with文はcontext managerに開始・終了処理を任せ、ファイルなどの資源を確実に解放する。
            \end{itemize}
            \paragraph{上から順の処理}
            \begin{enumerate}[leftmargin=1.8em]
            \item with 文で path.open('r', encoding='utf-8', errors='ignore') を管理しながら処理する。
\item   max(0, sum((1 for \_ in f)) - 1) を返す。
            \end{enumerate}
            \paragraph{代表コード断片}
            \begin{lstlisting}[language=Python]
def csv_row_count(path: Path) -> int:
    with path.open("r", encoding="utf-8", errors="ignore") as f:
        return max(0, sum(1 for _ in f) - 1)
\end{lstlisting}

\subsection{L873 関数 summarize\_path}
            \begin{itemize}[leftmargin=1.6em]
              \item 行範囲: L873--L893
              \item 所属: module直下
              \item 定義: \path|summarize_path(path: Path) -> Dict[str, object]|
              \item docstring: 明示されていない。
              \item このブロックが直接呼ぶ主な関数・メソッド: csv\_row\_count, exists, fspath, join, len, lower, lstrip, read\_csv, str
              \item 戻り値の要点: summary / summary
              \item 主なローカル代入名: col, header, suffix, summary
              \item 読み取る主なインスタンス属性: 特になし。
              \item 更新する主なインスタンス属性: 特になし。
              \item 明示的に送出する例外: 明示的raiseはない。
              \item 制御構造の規模: 条件分岐 2、ループ 0、try 1
            \end{itemize}
            \paragraph{この定義を読むためのPython構文}
            \begin{itemize}[leftmargin=1.6em]
            \item `->`以後は戻り値の型注釈であり、実行時の値を自動変換する命令ではない。
\item 内包表記は、反復・条件・値生成を一つの式で記述してcollectionを作る。
\item try文は例外が起き得る処理と、異常時または終了時の経路を分ける。
            \end{itemize}
            \paragraph{上から順の処理}
            \begin{enumerate}[leftmargin=1.8em]
            \item suffix に path.suffix.lower() の結果を代入する。
\item summary に \{'path': os.fspath(path), 'exists': path.exists(), 'kind': suffix.lstrip('.'), 'rows': '', 'columns': '', 'column\_names': ''\} の結果を代入する。
\item 条件 not path.exists() を判定し、真なら内部処理を行う。
\item   summary を返す。
\item 条件 suffix == '.csv' を判定し、真なら内部処理を行う。
\item   例外処理を伴う try ブロックを実行する。
\item     header に pd.read\_csv(path, nrows=0) の結果を代入する。
\item     summary['rows'] に csv\_row\_count(path) の結果を代入する。
\item     summary['columns'] に len(header.columns) の結果を代入する。
\item     summary['column\_names'] に ', '.join((str(col) for col in header.columns[:12])) の結果を代入する。
\item     Exceptionを捕捉した場合:
\item     Pass 文を実行する。
\item summary を返す。
            \end{enumerate}
            \paragraph{代表コード断片}
            \begin{lstlisting}[language=Python]
def summarize_path(path: Path) -> Dict[str, object]:
    suffix = path.suffix.lower()
    summary = {
        "path": os.fspath(path),
        "exists": path.exists(),
        "kind": suffix.lstrip("."),
        "rows": "",
        "columns": "",
        "column_names": "",
    }
    if not path.exists():
        return summary
    if suffix == ".csv":
        try:
            header = pd.read_csv(path, nrows=0)
            summary["rows"] = csv_row_count(path)
            summary["columns"] = len(header.columns)
            summary["column_names"] = ", ".join(str(col) for col in header.columns[:12])
        except Exception:
            pass
    return summary
\end{lstlisting}

\subsection{L896 関数 dataframe\_to\_markdown}
            \begin{itemize}[leftmargin=1.6em]
              \item 行範囲: L896--L907
              \item 所属: module直下
              \item 定義: \path|dataframe_to_markdown(df: pd.DataFrame) -> str|
              \item docstring: 明示されていない。
              \item このブロックが直接呼ぶ主な関数・メソッド: append, iterrows, join, len, replace, str
              \item 戻り値の要点: '\textbackslash{}n'.join(lines) / '(none)'
              \item 主なローカル代入名: \_, col, cols, lines, row, vals
              \item 読み取る主なインスタンス属性: 特になし。
              \item 更新する主なインスタンス属性: 特になし。
              \item 明示的に送出する例外: 明示的raiseはない。
              \item 制御構造の規模: 条件分岐 1、ループ 1、try 0
            \end{itemize}
            \paragraph{この定義を読むためのPython構文}
            \begin{itemize}[leftmargin=1.6em]
            \item `->`以後は戻り値の型注釈であり、実行時の値を自動変換する命令ではない。
\item 内包表記は、反復・条件・値生成を一つの式で記述してcollectionを作る。
            \end{itemize}
            \paragraph{上から順の処理}
            \begin{enumerate}[leftmargin=1.8em]
            \item 条件 df.empty を判定し、真なら内部処理を行う。
\item   '(none)' を返す。
\item cols に [str(col) for col in df.columns] の結果を代入する。
\item lines に ['| ' + ' | '.join(cols) + ' |', '| ' + ' | '.join(['---'] * len(cols)) + ' |'] の結果を代入する。
\item df.iterrows() を順に走査し、各要素を (\_, row) に入れて処理する。
\item   vals に [str(row[col]).replace('\textbackslash{}n', ' ') for col in df.columns] の結果を代入する。
\item   lines.append(...) を実行する。
\item '\textbackslash{}n'.join(lines) を返す。
            \end{enumerate}
            \paragraph{代表コード断片}
            \begin{lstlisting}[language=Python]
def dataframe_to_markdown(df: pd.DataFrame) -> str:
    if df.empty:
        return "(none)"
    cols = [str(col) for col in df.columns]
    lines = [
        "| " + " | ".join(cols) + " |",
        "| " + " | ".join(["---"] * len(cols)) + " |",
    ]
    for _, row in df.iterrows():
        vals = [str(row[col]).replace("\n", " ") for col in df.columns]
        lines.append("| " + " | ".join(vals) + " |")
    return "\n".join(lines)
\end{lstlisting}

\subsection{L910 関数 flatten\_scalars}
            \begin{itemize}[leftmargin=1.6em]
              \item 行範囲: L910--L917
              \item 所属: module直下
              \item 定義: \path|flatten_scalars(prefix: str, payload, rows: List[Dict[str, object]]) -> None|
              \item docstring: 明示されていない。
              \item このブロックが直接呼ぶ主な関数・メソッド: append, flatten\_scalars, isinstance, items, str
              \item 戻り値の要点: 戻り値が無いか、return 文が明示されていない。
              \item 主なローカル代入名: key, next\_prefix, value
              \item 読み取る主なインスタンス属性: 特になし。
              \item 更新する主なインスタンス属性: 特になし。
              \item 明示的に送出する例外: 明示的raiseはない。
              \item 制御構造の規模: 条件分岐 2、ループ 1、try 0
            \end{itemize}
            \paragraph{この定義を読むためのPython構文}
            \begin{itemize}[leftmargin=1.6em]
            \item `->`以後は戻り値の型注釈であり、実行時の値を自動変換する命令ではない。
            \end{itemize}
            \paragraph{上から順の処理}
            \begin{enumerate}[leftmargin=1.8em]
            \item 条件 isinstance(payload, dict) を判定し、真なら内部処理を行う。
\item   payload.items() を順に走査し、各要素を (key, value) に入れて処理する。
\item     next\_prefix に f'\{prefix\}.\{key\}' if prefix else str(key) の結果を代入する。
\item     flatten\_scalars(...) を実行する。
\item    を返す。
\item 条件 isinstance(payload, (int, float, bool, str)) を判定し、真なら内部処理を行う。
\item   rows.append(...) を実行する。
            \end{enumerate}
            \paragraph{代表コード断片}
            \begin{lstlisting}[language=Python]
def flatten_scalars(prefix: str, payload, rows: List[Dict[str, object]]) -> None:
    if isinstance(payload, dict):
        for key, value in payload.items():
            next_prefix = f"{prefix}.{key}" if prefix else str(key)
            flatten_scalars(next_prefix, value, rows)
        return
    if isinstance(payload, (int, float, bool, str)):
        rows.append({"key": prefix, "value": payload})
\end{lstlisting}

\subsection{L920 関数 build\_current\_asset\_manifests}
            \begin{itemize}[leftmargin=1.6em]
              \item 行範囲: L920--L983
              \item 所属: module直下
              \item 定義: \path|build_current_asset_manifests(profile_yaml: Path, fit_summary: dict, output_dir: Path) -> Dict[str, Path]|
              \item docstring: 明示されていない。
              \item このブロックが直接呼ぶ主な関数・メソッド: DataFrame, append, dataframe\_to\_markdown, flatten\_scalars, get, isinstance, isoformat, items, join, now, read\_yaml, repo\_relative
              \item 戻り値の要点: \{'files\_csv': files\_csv, 'scalars\_csv': scalars\_csv, 'markdown': md\_path\}
              \item 主なローカル代入名: file\_rows, files\_csv, files\_df, info, local\_fit\_rows, local\_rows, md\_lines, md\_path, path, paths\_cfg, profile\_cfg, rel\_path, role, row, scalar\_rows, scalars\_csv, scalars\_df, section, section\_name
              \item 読み取る主なインスタンス属性: 特になし。
              \item 更新する主なインスタンス属性: 特になし。
              \item 明示的に送出する例外: 明示的raiseはない。
              \item 制御構造の規模: 条件分岐 0、ループ 4、try 0
            \end{itemize}
            \paragraph{この定義を読むためのPython構文}
            \begin{itemize}[leftmargin=1.6em]
            \item `->`以後は戻り値の型注釈であり、実行時の値を自動変換する命令ではない。
            \end{itemize}
            \paragraph{上から順の処理}
            \begin{enumerate}[leftmargin=1.8em]
            \item profile\_cfg に read\_yaml(profile\_yaml) の結果を代入する。
\item paths\_cfg に profile\_cfg.get('paths', \{\}) if isinstance(profile\_cfg, dict) else \{\} の結果を代入する。
\item file\_rows に [] の結果を代入する。
\item sorted(paths\_cfg.items()) を順に走査し、各要素を (role, rel\_path) に入れて処理する。
\item   path に (profile\_yaml.parent / str(rel\_path)).resolve() の結果を代入する。
\item   info に summarize\_path(path) の結果を代入する。
\item   info['role'] に role の結果を代入する。
\item   file\_rows.append(...) を実行する。
\item files\_df に pd.DataFrame(file\_rows) の結果を代入する。
\item files\_csv に output\_dir / 'current\_active\_files.csv' の結果を代入する。
\item files\_df.to\_csv(...) を実行する。
\item scalar\_rows に [] を代入する。
\item ('model', 'mpc', 'runtime', 'simulation', 'live', 'measurement') を順に走査し、各要素を section\_name に入れて処理する。
\item   section に profile\_cfg.get(section\_name, \{\}) の結果を代入する。
\item   local\_rows に [] を代入する。
\item   flatten\_scalars(...) を実行する。
\item   local\_rows を順に走査し、各要素を row に入れて処理する。
\item     scalar\_rows.append(...) を実行する。
\item local\_fit\_rows に [] を代入する。
\item flatten\_scalars(...) を実行する。
\item local\_fit\_rows を順に走査し、各要素を row に入れて処理する。
\item   scalar\_rows.append(...) を実行する。
\item scalars\_df に pd.DataFrame(scalar\_rows) の結果を代入する。
\item scalars\_csv に output\_dir / 'current\_scalar\_coefficients.csv' の結果を代入する。
\item scalars\_df.to\_csv(...) を実行する。
\item md\_lines に ['\# Current maps and coefficients', '', f'- profile: `\{repo\_relative(profile\_yaml)\}`', f"- generated\_at: `\{datetime.now().isoformat(timespec='seconds')\}`", '', '\#\# Active files', '', dataframe\_to\_markdown(files\_df), '', '\#\# Scalar coefficients', '', dataframe\_to\_markdown(scalars\_df), ''] の結果を代入する。
\item md\_path に output\_dir / 'current\_maps\_and\_coefficients.md' の結果を代入する。
\item md\_path.write\_text(...) を実行する。
\item \{'files\_csv': files\_csv, 'scalars\_csv': scalars\_csv, 'markdown': md\_path\} を返す。
            \end{enumerate}
            \paragraph{代表コード断片}
            \begin{lstlisting}[language=Python]
def build_current_asset_manifests(profile_yaml: Path, fit_summary: dict, output_dir: Path) -> Dict[str, Path]:
    profile_cfg = read_yaml(profile_yaml)
    paths_cfg = profile_cfg.get("paths", {}) if isinstance(profile_cfg, dict) else {}

    file_rows = []
    for role, rel_path in sorted(paths_cfg.items()):
        path = (profile_yaml.parent / str(rel_path)).resolve()
        info = summarize_path(path)
        info["role"] = role
        file_rows.append(info)
    files_df = pd.DataFrame(file_rows)
    files_csv = output_dir / "current_active_files.csv"
    files_df.to_csv(files_csv, index=False)

    scalar_rows: List[Dict[str, object]] = []
    for section_name in ("model", "mpc", "runtime", "simulation", "live", "measurement"):
        section = profile_cfg.get(section_name, {})
        local_rows: List[Dict[str, object]] = []
        flatten_scalars("", section, local_rows)
        for row in local_rows:
            scalar_rows.append(
                {
                    "source": f"profile:{section_name}",
                    "key": row["key"],
                    "value": row["value"],
                }
            )
    local_fit_rows: List[Dict[str, object]] = []
    flatten_scalars("", fit_summary, local_fit_rows)
    for row in local_fit_rows:
        scalar_rows.append(
            {
                "source": "fit_summary",
                "key": row["key"],
                "value": row["value"],
...
\end{lstlisting}

\subsection{L986 関数 render\_report}
            \begin{itemize}[leftmargin=1.6em]
              \item 行範囲: L986--L1501
              \item 所属: module直下
              \item 定義: \path!render_report(output_dir: Path, source_profile_yaml: Path, eval_profile_yaml: Path, fit_summary: dict, baseline_result: Dict[str, object], tuned_result: Dict[str, object], trials_df: pd.DataFrame, best_weights: Dict[str, float], manifest_paths: Dict[str, Path], removed_reference: str, specs: List[TermSpec], search_cfg: Dict[str, object], tensorboard_dir: Path | None) -> tuple[Path, Path]!
              \item docstring: 明示されていない。
              \item このブロックが直接呼ぶ主な関数・メソッド: DataFrame, Path, Series, UpperCostConfig, abs, active\_upper\_cost\_terms, append, arange, as\_posix, bar, close, compile\_tex
              \item 戻り値の要点: (tex\_path, md\_path) / latex\_escape(os.path.relpath(path, report\_dir)).replace('\%', '\textbackslash{}\textbackslash{}\%')
              \item 主なローカル代入名: \_, active\_terms, actual\_retire\_km, base\_row, base\_scenarios, baseline\_series, best\_weights\_csv, best\_weights\_csv\_label, best\_weights\_df, estimated\_10000\_gen\_candidates, estimated\_10000\_gen\_cpu\_days, eval\_profile\_label, files\_csv\_label, fit\_csv, fit\_df, fit\_rows, human\_gap\_km, inactive\_terms, item, iter\_compare\_csv
              \item 読み取る主なインスタンス属性: 特になし。
              \item 更新する主なインスタンス属性: 特になし。
              \item 明示的に送出する例外: 明示的raiseはない。
              \item 制御構造の規模: 条件分岐 7、ループ 4、try 1
            \end{itemize}
            \paragraph{この定義を読むためのPython構文}
            \begin{itemize}[leftmargin=1.6em]
            \item `->`以後は戻り値の型注釈であり、実行時の値を自動変換する命令ではない。
\item 内包表記は、反復・条件・値生成を一つの式で記述してcollectionを作る。
\item try文は例外が起き得る処理と、異常時または終了時の経路を分ける。
            \end{itemize}
            \paragraph{上から順の処理}
            \begin{enumerate}[leftmargin=1.8em]
            \item score\_label に 'robust score' if str(search\_cfg.get('scenario\_mode', 'nominal')) == 'robust' else 'aggregate score' の結果を代入する。
\item report\_dir に output\_dir / 'report' の結果を代入する。
\item ensure\_dir(...) を実行する。
\item source\_profile\_label に repo\_relative(source\_profile\_yaml) の結果を代入する。
\item eval\_profile\_label に repo\_relative(eval\_profile\_yaml) の結果を代入する。
\item removed\_reference\_label に repo\_relative(removed\_reference) or '(none)' の結果を代入する。
\item files\_csv\_label に repo\_relative(manifest\_paths['files\_csv']) の結果を代入する。
\item scalars\_csv\_label に repo\_relative(manifest\_paths['scalars\_csv']) の結果を代入する。
\item markdown\_label に repo\_relative(manifest\_paths['markdown']) の結果を代入する。
\item best\_weights\_csv\_label に repo\_relative(output\_dir / 'best\_upper\_cost.csv') の結果を代入する。
\item tensorboard\_label に repo\_relative(tensorboard\_dir) if tensorboard\_dir else '' の結果を代入する。
\item actual\_retire\_km に float(fit\_summary.get('race\_distance', \{\}).get('actual\_retire\_km', 2831.0)) の結果を代入する。
\item power\_rmse\_fit に float(fit\_summary.get('validation\_metrics', \{\}).get('power\_rmse\_fit\_window\_w', fit\_summary.get('validation\_metrics', \{\}).get('power\_rmse\_clean\_w', float('nan')))) の結果を代入する。
\item voltage\_rmse\_fit に float(fit\_summary.get('validation\_metrics', \{\}).get('voltage\_rmse\_fit\_window\_v', fit\_summary.get('validation\_metrics', \{\}).get('voltage\_rmse\_clean\_v', float('nan')))) の結果を代入する。
\item median\_cpu\_sec に float(pd.to\_numeric(trials\_df.get('cpu\_sec', pd.Series(dtype=float)), errors='coerce').dropna().median()) if not trials\_df.empty else float('nan') の結果を代入する。
\item search\_candidates に int(sum((1 for \_, row in trials\_df.iterrows() if str(row.get('candidate', '')).startswith('g')))) の結果を代入する。
\item estimated\_10000\_gen\_candidates に int(search\_cfg.get('population', 0)) * 10000 の結果を代入する。
\item estimated\_10000\_gen\_cpu\_days に median\_cpu\_sec * estimated\_10000\_gen\_candidates / 86400.0 if math.isfinite(median\_cpu\_sec) and median\_cpu\_sec > 0.0 and (estimated\_10000\_gen\_candidates > 0) else float('nan') の結果を代入する。
\item human\_gap\_km に actual\_retire\_km - float(tuned\_result['final\_distance\_km']) の結果を代入する。
\item learning\_png に report\_dir / 'learning\_curve.png' の結果を代入する。
\item 条件 not trials\_df.empty を判定し、真なら内部処理を行う。
\item   plt.figure(...) を実行する。
\item   plt.plot(...) を実行する。
\item   plt.plot(...) を実行する。
\item   plt.xlabel(...) を実行する。
\item   plt.ylabel(...) を実行する。
\item   plt.title(...) を実行する。
\item   plt.grid(...) を実行する。
\item   plt.legend(...) を実行する。
\item   plt.tight\_layout(...) を実行する。
\item   plt.savefig(...) を実行する。
\item   plt.close(...) を実行する。
\item scenario\_compare\_png に report\_dir / 'scenario\_compare.png' の結果を代入する。
\item base\_scenarios に \{row['scenario']: row for row in baseline\_result['scenario\_results']\} の結果を代入する。
\item tuned\_scenarios に \{row['scenario']: row for row in tuned\_result['scenario\_results']\} の結果を代入する。
\item scenario\_names に list(base\_scenarios.keys()) の結果を代入する。
\item x に np.arange(len(scenario\_names)) の結果を代入する。
\item width に 0.35 の結果を代入する。
\item plt.figure(...) を実行する。
\item plt.bar(...) を実行する。
\item plt.bar(...) を実行する。
\item plt.xticks(...) を実行する。
\item plt.ylabel(...) を実行する。
\item plt.title(...) を実行する。
\item plt.grid(...) を実行する。
\item plt.legend(...) を実行する。
\item plt.tight\_layout(...) を実行する。
\item plt.savefig(...) を実行する。
\item plt.close(...) を実行する。
\item baseline\_series に pd.read\_csv(os.fspath(baseline\_result['nominal\_out\_csv'])) の結果を代入する。
\item tuned\_series に pd.read\_csv(os.fspath(tuned\_result['nominal\_out\_csv'])) の結果を代入する。
\item speed\_compare\_png に report\_dir / 'speed\_compare.png' の結果を代入する。
\item plt.figure(...) を実行する。
\item plt.plot(...) を実行する。
\item plt.plot(...) を実行する。
\item plt.xlabel(...) を実行する。
\item plt.ylabel(...) を実行する。
\item plt.title(...) を実行する。
\item plt.grid(...) を実行する。
\item plt.legend(...) を実行する。
\item plt.tight\_layout(...) を実行する。
\item plt.savefig(...) を実行する。
\item plt.close(...) を実行する。
\item soc\_compare\_png に report\_dir / 'soc\_compare.png' の結果を代入する。
\item plt.figure(...) を実行する。
\item plt.plot(...) を実行する。
\item plt.plot(...) を実行する。
\item plt.xlabel(...) を実行する。
\item plt.ylabel(...) を実行する。
\item plt.title(...) を実行する。
\item plt.grid(...) を実行する。
\item plt.legend(...) を実行する。
\item plt.tight\_layout(...) を実行する。
\item plt.savefig(...) を実行する。
\item plt.close(...) を実行する。
\item best\_weights\_df に pd.DataFrame(\{'term': list(best\_weights.keys()), 'value': [float(best\_weights[key]) for key in best\_weights]\}) の結果を代入する。
\item best\_weights\_csv に output\_dir / 'best\_upper\_cost.csv' の結果を代入する。
\item best\_weights\_df.to\_csv(...) を実行する。
\item scenario\_rows に [] の結果を代入する。
\item baseline\_result['scenario\_results'] を順に走査し、各要素を base\_row に入れて処理する。
            \end{enumerate}
            \paragraph{代表コード断片}
            \begin{lstlisting}[language=Python]
def render_report(
    output_dir: Path,
    source_profile_yaml: Path,
    eval_profile_yaml: Path,
    fit_summary: dict,
    baseline_result: Dict[str, object],
    tuned_result: Dict[str, object],
    trials_df: pd.DataFrame,
    best_weights: Dict[str, float],
    manifest_paths: Dict[str, Path],
    removed_reference: str,
    specs: List[TermSpec],
    search_cfg: Dict[str, object],
    tensorboard_dir: Path | None,
) -> tuple[Path, Path]:
    score_label = "robust score" if str(search_cfg.get("scenario_mode", "nominal")) == "robust" else "aggregate score"
    report_dir = output_dir / "report"
    ensure_dir(report_dir)
    source_profile_label = repo_relative(source_profile_yaml)
    eval_profile_label = repo_relative(eval_profile_yaml)
    removed_reference_label = repo_relative(removed_reference) or "(none)"
    files_csv_label = repo_relative(manifest_paths["files_csv"])
    scalars_csv_label = repo_relative(manifest_paths["scalars_csv"])
    markdown_label = repo_relative(manifest_paths["markdown"])
    best_weights_csv_label = repo_relative(output_dir / "best_upper_cost.csv")
    tensorboard_label = repo_relative(tensorboard_dir) if tensorboard_dir else ""
    actual_retire_km = float(fit_summary.get("race_distance", {}).get("actual_retire_km", 2831.0))
    power_rmse_fit = float(
        fit_summary.get("validation_metrics", {}).get(
            "power_rmse_fit_window_w",
            fit_summary.get("validation_metrics", {}).get("power_rmse_clean_w", float("nan")),
        )
    )
    voltage_rmse_fit = float(
        fit_summary.get("validation_metrics", {}).get(
...
\end{lstlisting}

\subsection{L1138 関数 render\_report.tex\_path\_rel}
            \begin{itemize}[leftmargin=1.6em]
              \item 行範囲: L1138--L1139
              \item 所属: \path|render_report|
              \item 定義: \path|tex_path_rel(path: Path) -> str|
              \item docstring: 明示されていない。
              \item このブロックが直接呼ぶ主な関数・メソッド: latex\_escape, relpath, replace
              \item 戻り値の要点: latex\_escape(os.path.relpath(path, report\_dir)).replace('\%', '\textbackslash{}\textbackslash{}\%')
              \item 主なローカル代入名: 特になし。
              \item 読み取る主なインスタンス属性: 特になし。
              \item 更新する主なインスタンス属性: 特になし。
              \item 明示的に送出する例外: 明示的raiseはない。
              \item 制御構造の規模: 条件分岐 0、ループ 0、try 0
            \end{itemize}
            \paragraph{この定義を読むためのPython構文}
            \begin{itemize}[leftmargin=1.6em]
            \item `->`以後は戻り値の型注釈であり、実行時の値を自動変換する命令ではない。
            \end{itemize}
            \paragraph{上から順の処理}
            \begin{enumerate}[leftmargin=1.8em]
            \item latex\_escape(os.path.relpath(path, report\_dir)).replace('\%', '\textbackslash{}\textbackslash{}\%') を返す。
            \end{enumerate}
            \paragraph{代表コード断片}
            \begin{lstlisting}[language=Python]
    def tex_path_rel(path: Path) -> str:
        return latex_escape(os.path.relpath(path, report_dir)).replace("%", r"\%")
\end{lstlisting}

\subsection{L1504 関数 main}
            \begin{itemize}[leftmargin=1.6em]
              \item 行範囲: L1504--L1937
              \item 所属: module直下
              \item 定義: \path|main() -> None|
              \item docstring: 明示されていない。
              \item このブロックが直接呼ぶ主な関数・メソッド: ArgumentParser, DataFrame, Path, RuntimeError, ScenarioRiskConfig, SummaryWriter, UpperCostConfig, add\_argument, append, array, bool, build\_current\_asset\_manifests
              \item 戻り値の要点: 戻り値が無いか、return 文が明示されていない。
              \item 主なローカル代入名: ap, args, backup\_profile, base\_cost\_cfg, base\_cost\_seed, base\_weights, baseline\_exact, best\_result, best\_weights\_yaml, coarse\_candidate, coarse\_cfg\_overrides, coarse\_cli\_overrides, coarse\_result, coarse\_trials, correlated\_grid, dotted\_key, elite, elite\_vecs, eval\_profile\_yaml, fit\_summary
              \item 読み取る主なインスタンス属性: 特になし。
              \item 更新する主なインスタンス属性: 特になし。
              \item 明示的に送出する例外: RuntimeError('No candidate was evaluated.'), RuntimeError('Validation stage produced no result.')
              \item 制御構造の規模: 条件分岐 15、ループ 6、try 0
            \end{itemize}
            \paragraph{この定義を読むためのPython構文}
            \begin{itemize}[leftmargin=1.6em]
            \item `->`以後は戻り値の型注釈であり、実行時の値を自動変換する命令ではない。
\item lambdaは名前を付けずに短い関数オブジェクトを作る構文である。
\item 内包表記は、反復・条件・値生成を一つの式で記述してcollectionを作る。
            \end{itemize}
            \paragraph{上から順の処理}
            \begin{enumerate}[leftmargin=1.8em]
            \item ap に argparse.ArgumentParser() の結果を代入する。
\item ap.add\_argument(...) を実行する。
\item ap.add\_argument(...) を実行する。
\item ap.add\_argument(...) を実行する。
\item ap.add\_argument(...) を実行する。
\item ap.add\_argument(...) を実行する。
\item ap.add\_argument(...) を実行する。
\item ap.add\_argument(...) を実行する。
\item ap.add\_argument(...) を実行する。
\item ap.add\_argument(...) を実行する。
\item ap.add\_argument(...) を実行する。
\item ap.add\_argument(...) を実行する。
\item ap.add\_argument(...) を実行する。
\item ap.add\_argument(...) を実行する。
\item ap.add\_argument(...) を実行する。
\item ap.add\_argument(...) を実行する。
\item ap.add\_argument(...) を実行する。
\item ap.add\_argument(...) を実行する。
\item ap.add\_argument(...) を実行する。
\item ap.add\_argument(...) を実行する。
\item args に ap.parse\_args() の結果を代入する。
\item profile\_yaml に Path(args.profile\_yaml).resolve() の結果を代入する。
\item profile\_cfg に read\_yaml(profile\_yaml) の結果を代入する。
\item base\_cost\_cfg に load\_upper\_cost\_config(profile\_cfg.get('mpc', \{\}), legacy=profile\_cfg.get('mpc', \{\})) の結果を代入する。
\item scenario\_mode に str(args.scenario\_mode).strip().lower() の結果を代入する。
\item risk\_config に ScenarioRiskConfig(cvar\_alpha=float(args.risk\_cvar\_alpha), cvar\_weight=float(args.risk\_cvar\_weight), variance\_weight=float(args.risk\_variance\_weight), max\_failure\_probability=float(args.max\_scenario\_failure\_probability), infeasible\_score=float(args.infeasible\_score)) の結果を代入する。
\item include\_terminal\_term に float(profile\_cfg.get('mpc', \{\}).get('soc\_finish\_target', -1.0)) > 0.0 の結果を代入する。
\item include\_uncertainty\_term に scenario\_mode != 'nominal' の結果を代入する。
\item base\_cost\_seed に base\_cost\_cfg.to\_dict() の結果を代入する。
\item 条件 not include\_uncertainty\_term を判定し、真なら内部処理を行う。
\item   base\_cost\_seed['w\_uncertainty\_reserve'] に 0.0 の結果を代入する。
\item   base\_cost\_seed['reserve\_soc\_per\_hour'] に 0.0 の結果を代入する。
\item   base\_cost\_seed['reserve\_soc\_max\_extra'] に 0.0 の結果を代入する。
\item 条件 not include\_terminal\_term を判定し、真なら内部処理を行う。
\item   base\_cost\_seed['w\_soc\_terminal'] に 0.0 の結果を代入する。
\item base\_cost\_cfg に UpperCostConfig(**base\_cost\_seed) の結果を代入する。
\item specs に upper\_cost\_specs(base\_cost\_cfg, include\_progress\_terms=args.include\_progress\_terms, include\_uncertainty\_term=include\_uncertainty\_term, include\_terminal\_term=include\_terminal\_term) の結果を代入する。
\item rng に np.random.default\_rng(args.seed) の結果を代入する。
\item planning\_race\_km に float(profile\_cfg.get('mpc', \{\}).get('race\_km', 3035.5)) の結果を代入する。
\item package\_dir に profile\_yaml.parent の結果を代入する。
\item timestamp に datetime.now().strftime('\%Y\%m\%d\_\%H\%M\%S') の結果を代入する。
\item output\_dir に package\_dir / 'outputs' / 'self\_learning\_upper' / timestamp の結果を代入する。
\item ensure\_dir(...) を実行する。
\item tensorboard\_dir に output\_dir / 'tensorboard' の結果を代入する。
\item tb\_writer に SummaryWriter(log\_dir=os.fspath(tensorboard\_dir)) if SummaryWriter is not None else None の結果を代入する。
\item (eval\_profile\_yaml, removed\_reference) に build\_reference\_free\_profile(profile\_yaml, output\_dir, disable\_uncertainty\_reserve=scenario\_mode == 'nominal') の結果を代入する。
\item fit\_summary\_path に package\_dir / 'outputs' / 'identification' / f'\{package\_dir.name\}\_fit\_summary.yaml' の結果を代入する。
\item fit\_summary に read\_yaml(fit\_summary\_path) if fit\_summary\_path.exists() else \{\} の結果を代入する。
\item manifest\_paths に build\_current\_asset\_manifests(profile\_yaml, fit\_summary, output\_dir) の結果を代入する。
\item coarse\_cfg\_overrides に \{'mpc.dt': 5400.0, 'mpc.upper\_horizon\_mode': 'adaptive\_full\_race', 'mpc.upper\_max\_iter': int(args.coarse\_upper\_max\_iter), 'mpc.upper\_max\_steps': 6, 'mpc.race\_km': planning\_race\_km, 'mpc.upper\_horizon\_km': planning\_race\_km, 'mpc.upper\_ctrl\_km': 700.0, 'mpc.upper\_replan\_km': 0.0, 'mpc.upper\_replan\_sec': 0.0, 'mpc.upper\_adaptive\_min\_ds\_km': 20.0, 'mpc.upper\_adaptive\_max\_ds\_km': 400.0, 'mpc.upper\_adaptive\_growth': 1.3, 'mpc.reference\_speed\_tracking.enabled': False, 'mpc.upper\_cost.w\_uncertainty\_reserve': 0.0 if scenario\_mode == 'nominal' else float(base\_cost\_cfg.w\_uncertainty\_reserve), 'mpc.upper\_cost.reserve\_soc\_per\_hour': 0.0 if scenario\_mode == 'nominal' else float(base\_cost\_cfg.reserve\_soc\_per\_hour), 'mpc.upper\_cost.reserve\_soc\_max\_extra': 0.0 if scenario\_mode == 'nominal' else float(base\_cost\_cfg.reserve\_soc\_max\_extra)\} の結果を代入する。
\item coarse\_cli\_overrides に \{\} の結果を代入する。
\item validation\_cfg\_overrides に \{'mpc.dt': 2400.0, 'mpc.upper\_horizon\_mode': 'adaptive\_full\_race', 'mpc.upper\_max\_iter': int(args.validation\_upper\_max\_iter), 'mpc.upper\_max\_steps': 10, 'mpc.race\_km': planning\_race\_km, 'mpc.upper\_horizon\_km': planning\_race\_km, 'mpc.upper\_ctrl\_km': 500.0, 'mpc.upper\_replan\_km': 0.0, 'mpc.upper\_replan\_sec': 0.0, 'mpc.upper\_adaptive\_min\_ds\_km': 20.0, 'mpc.upper\_adaptive\_max\_ds\_km': 350.0, 'mpc.upper\_adaptive\_growth': 1.28, 'mpc.reference\_speed\_tracking.enabled': False, 'mpc.upper\_cost.w\_uncertainty\_reserve': 0.0 if scenario\_mode == 'nominal' else float(base\_cost\_cfg.w\_uncertainty\_reserve), 'mpc.upper\_cost.reserve\_soc\_per\_hour': 0.0 if scenario\_mode == 'nominal' else float(base\_cost\_cfg.reserve\_soc\_per\_hour), 'mpc.upper\_cost.reserve\_soc\_max\_extra': 0.0 if scenario\_mode == 'nominal' else float(base\_cost\_cfg.reserve\_soc\_max\_extra)\} の結果を代入する。
\item validation\_cli\_overrides に \{\} の結果を代入する。
\item medium\_cfg\_overrides に \{'mpc.dt': 3600.0, 'mpc.upper\_horizon\_mode': 'adaptive\_full\_race', 'mpc.upper\_max\_iter': int(args.elite\_medium\_upper\_max\_iter), 'mpc.upper\_max\_steps': 8, 'mpc.race\_km': planning\_race\_km, 'mpc.upper\_horizon\_km': planning\_race\_km, 'mpc.upper\_ctrl\_km': 600.0, 'mpc.upper\_replan\_km': 0.0, 'mpc.upper\_replan\_sec': 0.0, 'mpc.upper\_adaptive\_min\_ds\_km': 20.0, 'mpc.upper\_adaptive\_max\_ds\_km': 380.0, 'mpc.upper\_adaptive\_growth': 1.29, 'mpc.reference\_speed\_tracking.enabled': False, 'mpc.upper\_cost.w\_uncertainty\_reserve': 0.0 if scenario\_mode == 'nominal' else float(base\_cost\_cfg.w\_uncertainty\_reserve), 'mpc.upper\_cost.reserve\_soc\_per\_hour': 0.0 if scenario\_mode == 'nominal' else float(base\_cost\_cfg.reserve\_soc\_per\_hour), 'mpc.upper\_cost.reserve\_soc\_max\_extra': 0.0 if scenario\_mode == 'nominal' else float(base\_cost\_cfg.reserve\_soc\_max\_extra)\} の結果を代入する。
\item medium\_cli\_overrides に \{\} の結果を代入する。
\item 条件 args.fidelity\_mode == 'correlated' を判定し、真なら内部処理を行う。
\item   correlated\_grid に \{'mpc.dt': 2400.0, 'mpc.upper\_max\_steps': 10, 'mpc.upper\_ctrl\_km': 500.0, 'mpc.upper\_adaptive\_min\_ds\_km': 20.0, 'mpc.upper\_adaptive\_max\_ds\_km': 350.0, 'mpc.upper\_adaptive\_growth': 1.28\} の結果を代入する。
\item   coarse\_cfg\_overrides.update(...) を実行する。
\item   medium\_cfg\_overrides.update(...) を実行する。
\item   validation\_cfg\_overrides.update(...) を実行する。
\item scenarios に default\_scenarios(eval\_profile\_yaml, mode=scenario\_mode) の結果を代入する。
\item base\_weights に base\_cost\_cfg.to\_dict() の結果を代入する。
\item 条件 scenario\_mode == 'nominal' を判定し、真なら内部処理を行う。
\item   base\_weights['w\_uncertainty\_reserve'] に 0.0 の結果を代入する。
\item   base\_weights['reserve\_soc\_per\_hour'] に 0.0 の結果を代入する。
\item   base\_weights['reserve\_soc\_max\_extra'] に 0.0 の結果を代入する。
\item baseline\_exact に run\_candidate(eval\_profile\_yaml, output\_dir, 'validation\_baseline', base\_weights, validation\_cfg\_overrides, validation\_cli\_overrides, scenarios, risk\_config) の結果を代入する。
\item baseline\_exact['generation'] に -1 の結果を代入する。
\item baseline\_exact['trial\_index'] に -1 の結果を代入する。
\item log\_trial\_to\_tensorboard(...) を実行する。
\item mean に np.array([spec.init\_log10 for spec in specs], dtype=float) の結果を代入する。
\item sigma に np.array([0.75] * len(specs), dtype=float) の結果を代入する。
\item trials に [] の結果を代入する。
\item best\_result に None の結果を代入する。
\item trial\_index に 0 の結果を代入する。
\item range(args.generations) を順に走査し、各要素を generation に入れて処理する。
\item   generation\_results に [] の結果を代入する。
\item   range(args.population) を順に走査し、各要素を pop\_idx に入れて処理する。
\item     条件 generation == 0 and pop\_idx == 0 を判定し、真なら内部処理を行う。
\item       vec に mean.copy() の結果を代入する。
            \end{enumerate}
            \paragraph{代表コード断片}
            \begin{lstlisting}[language=Python]
def main() -> None:
    ap = argparse.ArgumentParser()
    ap.add_argument("--profile_yaml", default=os.fspath(DEFAULT_PROFILE))
    ap.add_argument("--output_profile_yaml", default="")
    ap.add_argument("--generations", type=int, default=16)
    ap.add_argument("--population", type=int, default=8)
    ap.add_argument("--elite_count", type=int, default=3)
    ap.add_argument("--validation_top_k", type=int, default=5)
    ap.add_argument("--elite_medium_top_k", type=int, default=4)
    ap.add_argument("--seed", type=int, default=7)
    ap.add_argument("--include_progress_terms", action="store_true")
    ap.add_argument("--scenario-mode", choices=["nominal", "robust"], default="nominal")
    ap.add_argument("--coarse-upper-max-iter", type=int, default=1)
    ap.add_argument("--elite-medium-upper-max-iter", type=int, default=3)
    ap.add_argument("--validation-upper-max-iter", type=int, default=8)
    ap.add_argument("--risk-cvar-alpha", type=float, default=0.90)
    ap.add_argument("--risk-cvar-weight", type=float, default=0.30)
    ap.add_argument("--risk-variance-weight", type=float, default=0.0)
    ap.add_argument("--max-scenario-failure-probability", type=float, default=0.01)
    ap.add_argument("--infeasible-score", type=float, default=-1.0e9)
    ap.add_argument(
        "--fidelity-mode",
        choices=["correlated", "legacy_fast"],
        default="correlated",
        help="Use identical discretization across fidelities so candidate ranking remains valid.",
    )
    args = ap.parse_args()

    profile_yaml = Path(args.profile_yaml).resolve()
    profile_cfg = read_yaml(profile_yaml)
    base_cost_cfg = load_upper_cost_config(profile_cfg.get("mpc", {}), legacy=profile_cfg.get("mpc", {}))
    scenario_mode = str(args.scenario_mode).strip().lower()
    risk_config = ScenarioRiskConfig(
        cvar_alpha=float(args.risk_cvar_alpha),
        cvar_weight=float(args.risk_cvar_weight),
...
\end{lstlisting}







        \subsection{CLI 引数}
            \begin{longtable}{p{0.07\linewidth} p{0.78\linewidth}}
            \toprule
            行 & 引数 \\
            \midrule
            1506 & \path|--profile_yaml| \\
1507 & \path|--output_profile_yaml| \\
1508 & \path|--generations| \\
1509 & \path|--population| \\
1510 & \path|--elite_count| \\
1511 & \path|--validation_top_k| \\
1512 & \path|--elite_medium_top_k| \\
1513 & \path|--seed| \\
1514 & \path|--include_progress_terms| \\
1515 & \path|--scenario-mode| \\
1516 & \path|--coarse-upper-max-iter| \\
1517 & \path|--elite-medium-upper-max-iter| \\
1518 & \path|--validation-upper-max-iter| \\
1519 & \path|--risk-cvar-alpha| \\
1520 & \path|--risk-cvar-weight| \\
1521 & \path|--risk-variance-weight| \\
1522 & \path|--max-scenario-failure-probability| \\
1523 & \path|--infeasible-score| \\
1524 & \path|--fidelity-mode| \\
            \bottomrule
            \end{longtable}



        \section{処理の流れ}
        \begin{enumerate}[leftmargin=1.7em]
        \item CLI 引数を解釈し、main() から処理を起動する。
        \end{enumerate}

        \section{このファイルの位置づけ}
        この資料群では、\path|scripts/tune_upper_planner_weights.py| を
        \textbf{planning research}の中核として扱う。
        したがって、ここで扱う処理は単独で完結せず、
        前段の profile / launch / telemetry と後段の planner / logger / report へ繋がる。

        \end{document}
